\documentclass[12pt,a4paper]{report}
\usepackage[utf8]{inputenc}
\usepackage[margin=1in]{geometry}
\usepackage{amsmath,amssymb,amsthm}
\usepackage{graphicx}
\usepackage{hyperref}
\usepackage{enumitem}
\usepackage{tcolorbox}
\usepackage{fancyhdr}
\usepackage{tikz}
\usepackage{algorithm}
\usepackage{algpseudocode}
\usepackage{listings}
\usepackage{xcolor}

% Page style
\pagestyle{fancy}
\fancyhf{}
\rhead{GATE CSE Study Guide}
\lhead{\leftmark}
\cfoot{\thepage}

% Custom environments
\newtcolorbox{concept}{colback=blue!5!white,colframe=blue!75!black,title=Key Concept}
\newtcolorbox{theorem}{colback=green!5!white,colframe=green!75!black,title=Theorem}
\newtcolorbox{formula}{colback=red!5!white,colframe=red!75!black,title=Important Formula}

% Title page
\title{\Huge\textbf{GATE CSE}\\\LARGE Complete Theoretical Concepts\\Study Guide}
\author{Comprehensive Reference for Computer Science Engineering}
\date{\today}

\begin{document}

\maketitle

\tableofcontents

\chapter{Engineering Mathematics}

\section{Discrete Mathematics}

\subsection{Propositional Logic}
\begin{concept}
\textbf{Propositional Logic} deals with propositions that can be either true or false.
\begin{itemize}
    \item \textbf{Proposition}: A declarative statement that is either true or false
    \item \textbf{Logical Connectives}: AND ($\land$), OR ($\lor$), NOT ($\neg$), IMPLIES ($\rightarrow$), IFF ($\leftrightarrow$)
    \item \textbf{Tautology}: A proposition that is always true
    \item \textbf{Contradiction}: A proposition that is always false
    \item \textbf{Contingency}: A proposition that can be true or false depending on values
\end{itemize}
\end{concept}

\textbf{Laws of Logic:}
\begin{itemize}
    \item Identity Laws: $p \land T \equiv p$, $p \lor F \equiv p$
    \item Domination Laws: $p \lor T \equiv T$, $p \land F \equiv F$
    \item Idempotent Laws: $p \lor p \equiv p$, $p \land p \equiv p$
    \item Double Negation: $\neg(\neg p) \equiv p$
    \item Commutative Laws: $p \lor q \equiv q \lor p$, $p \land q \equiv q \land p$
    \item Associative Laws: $(p \lor q) \lor r \equiv p \lor (q \lor r)$
    \item Distributive Laws: $p \lor (q \land r) \equiv (p \lor q) \land (p \lor r)$
    \item De Morgan's Laws: $\neg(p \land q) \equiv \neg p \lor \neg q$, $\neg(p \lor q) \equiv \neg p \land \neg q$
\end{itemize}

\subsection{Predicate Logic}
\begin{concept}
\textbf{Predicate Logic} extends propositional logic with predicates and quantifiers.
\begin{itemize}
    \item \textbf{Universal Quantifier} ($\forall$): "for all"
    \item \textbf{Existential Quantifier} ($\exists$): "there exists"
    \item \textbf{Negation Rules}: $\neg \forall x P(x) \equiv \exists x \neg P(x)$
    \item $\neg \exists x P(x) \equiv \forall x \neg P(x)$
\end{itemize}
\end{concept}

\subsection{Set Theory}
\begin{concept}
\textbf{Sets} are collections of distinct objects.
\begin{itemize}
    \item \textbf{Operations}: Union ($\cup$), Intersection ($\cap$), Difference ($-$), Complement ($'$)
    \item \textbf{Properties}: $|A \cup B| = |A| + |B| - |A \cap B|$
    \item \textbf{Power Set}: Set of all subsets, $|P(A)| = 2^{|A|}$
    \item \textbf{Cartesian Product}: $A \times B = \{(a,b) : a \in A, b \in B\}$
\end{itemize}
\end{concept}

\subsection{Relations}
\begin{concept}
A \textbf{Relation} R from set A to set B is a subset of $A \times B$.
\begin{itemize}
    \item \textbf{Reflexive}: $\forall x \in A: (x,x) \in R$
    \item \textbf{Symmetric}: $(x,y) \in R \Rightarrow (y,x) \in R$
    \item \textbf{Transitive}: $(x,y) \in R \land (y,z) \in R \Rightarrow (x,z) \in R$
    \item \textbf{Anti-symmetric}: $(x,y) \in R \land (y,x) \in R \Rightarrow x = y$
    \item \textbf{Equivalence Relation}: Reflexive, Symmetric, and Transitive
    \item \textbf{Partial Order}: Reflexive, Anti-symmetric, and Transitive
\end{itemize}
\end{concept}

\subsection{Functions}
\begin{concept}
Types of \textbf{Functions}:
\begin{itemize}
    \item \textbf{One-to-One (Injective)}: $f(x_1) = f(x_2) \Rightarrow x_1 = x_2$
    \item \textbf{Onto (Surjective)}: $\forall y \in B, \exists x \in A: f(x) = y$
    \item \textbf{Bijective}: Both one-to-one and onto
    \item \textbf{Inverse}: Exists only for bijective functions
\end{itemize}
\end{concept}

\subsection{Combinatorics}
\begin{formula}
\textbf{Basic Counting Principles:}
\begin{itemize}
    \item \textbf{Permutations}: $P(n,r) = \frac{n!}{(n-r)!}$
    \item \textbf{Combinations}: $C(n,r) = \binom{n}{r} = \frac{n!}{r!(n-r)!}$
    \item \textbf{Pigeonhole Principle}: If $n$ items are placed in $m$ containers and $n > m$, at least one container has more than one item
    \item \textbf{Inclusion-Exclusion}: $|A \cup B| = |A| + |B| - |A \cap B|$
\end{itemize}
\end{formula}

\subsection{Graph Theory}
\begin{concept}
A \textbf{Graph} $G = (V, E)$ consists of vertices and edges.
\begin{itemize}
    \item \textbf{Degree}: Number of edges incident on a vertex
    \item \textbf{Handshaking Lemma}: $\sum_{v \in V} deg(v) = 2|E|$
    \item \textbf{Path}: Sequence of vertices connected by edges
    \item \textbf{Cycle}: Path that starts and ends at the same vertex
    \item \textbf{Connected Graph}: Path exists between every pair of vertices
    \item \textbf{Tree}: Connected acyclic graph with $|E| = |V| - 1$
    \item \textbf{Spanning Tree}: Subgraph that includes all vertices and is a tree
    \item \textbf{Euler Path}: Path using every edge exactly once
    \item \textbf{Euler Circuit}: Euler path that starts and ends at same vertex (exists if all vertices have even degree)
    \item \textbf{Hamiltonian Path}: Path visiting every vertex exactly once
\end{itemize}
\end{concept}

\section{Linear Algebra}

\subsection{Matrices}
\begin{concept}
\textbf{Matrix Operations:}
\begin{itemize}
    \item \textbf{Addition}: $(A + B)_{ij} = A_{ij} + B_{ij}$
    \item \textbf{Multiplication}: $(AB)_{ij} = \sum_k A_{ik}B_{kj}$
    \item \textbf{Transpose}: $(A^T)_{ij} = A_{ji}$
    \item \textbf{Determinant}: Scalar value, $det(AB) = det(A) \cdot det(B)$
    \item \textbf{Inverse}: $AA^{-1} = I$ (exists if $det(A) \neq 0$)
\end{itemize}
\end{concept}

\begin{concept}
\textbf{Special Matrices:}
\begin{itemize}
    \item \textbf{Identity Matrix}: $I_{ij} = 1$ if $i=j$, else 0
    \item \textbf{Symmetric}: $A = A^T$
    \item \textbf{Skew-Symmetric}: $A = -A^T$
    \item \textbf{Orthogonal}: $AA^T = I$
    \item \textbf{Diagonal}: Non-zero elements only on diagonal
\end{itemize}
\end{concept}

\subsection{Vector Spaces}
\begin{concept}
A \textbf{Vector Space} is a set with addition and scalar multiplication operations.
\begin{itemize}
    \item \textbf{Linear Independence}: Vectors that cannot be expressed as linear combinations of others
    \item \textbf{Basis}: Minimal set of linearly independent vectors spanning the space
    \item \textbf{Dimension}: Number of vectors in a basis
    \item \textbf{Rank}: Dimension of column space (or row space)
    \item \textbf{Nullity}: Dimension of null space
\end{itemize}
\end{concept}

\begin{theorem}
\textbf{Rank-Nullity Theorem}: For matrix $A: \mathbb{R}^n \to \mathbb{R}^m$
$$rank(A) + nullity(A) = n$$
\end{theorem}

\subsection{Eigenvalues and Eigenvectors}
\begin{concept}
For matrix $A$, if $Av = \lambda v$ (where $v \neq 0$):
\begin{itemize}
    \item $\lambda$ is an \textbf{eigenvalue}
    \item $v$ is an \textbf{eigenvector}
    \item \textbf{Characteristic Equation}: $det(A - \lambda I) = 0$
    \item Sum of eigenvalues = Trace of matrix
    \item Product of eigenvalues = Determinant of matrix
\end{itemize}
\end{concept}

\section{Calculus}

\subsection{Limits and Continuity}
\begin{concept}
\begin{itemize}
    \item \textbf{Limit}: $\lim_{x \to a} f(x) = L$ if $f(x)$ approaches $L$ as $x$ approaches $a$
    \item \textbf{Continuity}: $f$ is continuous at $a$ if $\lim_{x \to a} f(x) = f(a)$
    \item \textbf{L'Hôpital's Rule}: For $\frac{0}{0}$ or $\frac{\infty}{\infty}$ forms: $\lim_{x \to a} \frac{f(x)}{g(x)} = \lim_{x \to a} \frac{f'(x)}{g'(x)}$
\end{itemize}
\end{concept}

\subsection{Differentiation}
\begin{formula}
\textbf{Key Derivatives:}
\begin{itemize}
    \item $\frac{d}{dx}(x^n) = nx^{n-1}$
    \item $\frac{d}{dx}(e^x) = e^x$
    \item $\frac{d}{dx}(\ln x) = \frac{1}{x}$
    \item $\frac{d}{dx}(\sin x) = \cos x$
    \item $\frac{d}{dx}(\cos x) = -\sin x$
\end{itemize}
\textbf{Rules:}
\begin{itemize}
    \item Product Rule: $(uv)' = u'v + uv'$
    \item Quotient Rule: $(\frac{u}{v})' = \frac{u'v - uv'}{v^2}$
    \item Chain Rule: $(f(g(x)))' = f'(g(x)) \cdot g'(x)$
\end{itemize}
\end{formula}

\subsection{Applications of Derivatives}
\begin{concept}
\begin{itemize}
    \item \textbf{Critical Points}: Where $f'(x) = 0$ or undefined
    \item \textbf{Maxima/Minima}: Use first or second derivative test
    \item \textbf{Mean Value Theorem}: If $f$ continuous on $[a,b]$ and differentiable on $(a,b)$, then $\exists c \in (a,b)$ such that $f'(c) = \frac{f(b)-f(a)}{b-a}$
\end{itemize}
\end{concept}

\subsection{Integration}
\begin{formula}
\textbf{Key Integrals:}
\begin{itemize}
    \item $\int x^n dx = \frac{x^{n+1}}{n+1} + C$ (for $n \neq -1$)
    \item $\int e^x dx = e^x + C$
    \item $\int \frac{1}{x} dx = \ln|x| + C$
    \item $\int \sin x dx = -\cos x + C$
    \item $\int \cos x dx = \sin x + C$
\end{itemize}
\end{formula}

\begin{theorem}
\textbf{Fundamental Theorem of Calculus}:
$$\int_a^b f(x)dx = F(b) - F(a)$$
where $F'(x) = f(x)$
\end{theorem}

\section{Probability and Statistics}

\subsection{Probability Basics}
\begin{concept}
\begin{itemize}
    \item \textbf{Sample Space}: Set of all possible outcomes
    \item \textbf{Event}: Subset of sample space
    \item \textbf{Probability}: $P(A) = \frac{\text{favorable outcomes}}{\text{total outcomes}}$
    \item \textbf{Complement}: $P(A') = 1 - P(A)$
    \item \textbf{Addition Rule}: $P(A \cup B) = P(A) + P(B) - P(A \cap B)$
    \item \textbf{Conditional Probability}: $P(A|B) = \frac{P(A \cap B)}{P(B)}$
    \item \textbf{Independent Events}: $P(A \cap B) = P(A) \cdot P(B)$
\end{itemize}
\end{concept}

\begin{theorem}
\textbf{Bayes' Theorem}:
$$P(A|B) = \frac{P(B|A) \cdot P(A)}{P(B)}$$
\end{theorem}

\subsection{Random Variables}
\begin{concept}
\begin{itemize}
    \item \textbf{Random Variable}: Function mapping outcomes to real numbers
    \item \textbf{Discrete RV}: Countable values
    \item \textbf{Continuous RV}: Uncountable values in an interval
    \item \textbf{Probability Mass Function (PMF)}: For discrete RV
    \item \textbf{Probability Density Function (PDF)}: For continuous RV
    \item \textbf{Cumulative Distribution Function (CDF)}: $F(x) = P(X \leq x)$
\end{itemize}
\end{concept}

\subsection{Expected Value and Variance}
\begin{formula}
\begin{itemize}
    \item \textbf{Expected Value}: $E[X] = \sum x \cdot P(X=x)$ (discrete) or $\int x \cdot f(x)dx$ (continuous)
    \item \textbf{Variance}: $Var(X) = E[(X - E[X])^2] = E[X^2] - (E[X])^2$
    \item \textbf{Standard Deviation}: $\sigma = \sqrt{Var(X)}$
    \item \textbf{Properties}: $E[aX + b] = aE[X] + b$, $Var(aX + b) = a^2Var(X)$
\end{itemize}
\end{formula}

\subsection{Common Distributions}
\begin{concept}
\textbf{Discrete Distributions:}
\begin{itemize}
    \item \textbf{Bernoulli}: Single trial, $P(X=1) = p$, $E[X] = p$, $Var(X) = p(1-p)$
    \item \textbf{Binomial}: $n$ independent Bernoulli trials, $P(X=k) = \binom{n}{k}p^k(1-p)^{n-k}$
    \item \textbf{Poisson}: Rate $\lambda$, $P(X=k) = \frac{e^{-\lambda}\lambda^k}{k!}$, $E[X] = Var(X) = \lambda$
\end{itemize}

\textbf{Continuous Distributions:}
\begin{itemize}
    \item \textbf{Uniform}: $f(x) = \frac{1}{b-a}$ for $x \in [a,b]$, $E[X] = \frac{a+b}{2}$
    \item \textbf{Exponential}: $f(x) = \lambda e^{-\lambda x}$, $E[X] = \frac{1}{\lambda}$
    \item \textbf{Normal/Gaussian}: $f(x) = \frac{1}{\sigma\sqrt{2\pi}}e^{-\frac{(x-\mu)^2}{2\sigma^2}}$, $E[X] = \mu$, $Var(X) = \sigma^2$
\end{itemize}
\end{concept}

\chapter{Digital Logic}

\section{Boolean Algebra}

\subsection{Basic Operations}
\begin{concept}
\textbf{Boolean Operations:}
\begin{itemize}
    \item \textbf{AND} (·): $A \cdot B = 1$ only if both $A=1$ and $B=1$
    \item \textbf{OR} (+): $A + B = 1$ if either $A=1$ or $B=1$
    \item \textbf{NOT} ('): $A' = 1$ if $A=0$
    \item \textbf{NAND}: $(A \cdot B)'$
    \item \textbf{NOR}: $(A + B)'$
    \item \textbf{XOR} ($\oplus$): $A \oplus B = A'B + AB'$
    \item \textbf{XNOR}: $(A \oplus B)'$
\end{itemize}
\end{concept}

\subsection{Boolean Laws}
\begin{formula}
\textbf{Fundamental Laws:}
\begin{itemize}
    \item Identity: $A + 0 = A$, $A \cdot 1 = A$
    \item Null: $A + 1 = 1$, $A \cdot 0 = 0$
    \item Idempotent: $A + A = A$, $A \cdot A = A$
    \item Complement: $A + A' = 1$, $A \cdot A' = 0$
    \item Involution: $(A')' = A$
    \item Commutative: $A + B = B + A$, $A \cdot B = B \cdot A$
    \item Associative: $(A + B) + C = A + (B + C)$
    \item Distributive: $A \cdot (B + C) = A \cdot B + A \cdot C$
    \item De Morgan's: $(A + B)' = A' \cdot B'$, $(A \cdot B)' = A' + B'$
    \item Absorption: $A + A \cdot B = A$, $A \cdot (A + B) = A$
\end{itemize}
\end{formula}

\subsection{Canonical Forms}
\begin{concept}
\begin{itemize}
    \item \textbf{Minterm}: Product term where all variables appear (e.g., $A'B'C$)
    \item \textbf{Maxterm}: Sum term where all variables appear (e.g., $A+B+C'$)
    \item \textbf{Sum of Products (SOP)}: OR of minterms
    \item \textbf{Product of Sums (POS)}: AND of maxterms
    \item \textbf{Canonical SOP}: SOP with all minterms
    \item \textbf{Canonical POS}: POS with all maxterms
\end{itemize}
\end{concept}

\section{Logic Gates}

\subsection{Basic Gates}
\begin{concept}
\textbf{Gate Characteristics:}
\begin{itemize}
    \item \textbf{Universal Gates}: NAND and NOR (can implement any Boolean function)
    \item \textbf{Propagation Delay}: Time for output to change after input changes
    \item \textbf{Fan-in}: Maximum number of inputs a gate can have
    \item \textbf{Fan-out}: Maximum number of gates an output can drive
\end{itemize}
\end{concept}

\section{Combinational Circuits}

\subsection{Arithmetic Circuits}
\begin{concept}
\textbf{Half Adder}: Adds two bits
\begin{itemize}
    \item Sum: $S = A \oplus B$
    \item Carry: $C = A \cdot B$
\end{itemize}

\textbf{Full Adder}: Adds three bits (including carry-in)
\begin{itemize}
    \item Sum: $S = A \oplus B \oplus C_{in}$
    \item Carry: $C_{out} = AB + BC_{in} + AC_{in}$
\end{itemize}

\textbf{Ripple Carry Adder}: Cascade of full adders (delay = $n \times t_{FA}$)

\textbf{Carry Look-Ahead Adder}: Faster addition using parallel carry computation
\end{concept}

\subsection{Multiplexers and Demultiplexers}
\begin{concept}
\textbf{Multiplexer (MUX)}:
\begin{itemize}
    \item Selects one of $2^n$ inputs using $n$ select lines
    \item $2^n$-to-1 MUX requires $2^n$ data inputs and $n$ select lines
    \item Can implement any Boolean function
\end{itemize}

\textbf{Demultiplexer (DEMUX)}:
\begin{itemize}
    \item Routes one input to one of $2^n$ outputs
    \item 1-to-$2^n$ DEMUX requires 1 data input and $n$ select lines
\end{itemize}
\end{concept}

\subsection{Encoders and Decoders}
\begin{concept}
\textbf{Encoder}:
\begin{itemize}
    \item Converts $2^n$ inputs to $n$ outputs
    \item Priority Encoder: Encodes highest priority active input
\end{itemize}

\textbf{Decoder}:
\begin{itemize}
    \item Converts $n$ inputs to $2^n$ outputs
    \item Exactly one output is high for each input combination
    \item $n$-to-$2^n$ decoder activates one of $2^n$ output lines
\end{itemize}
\end{concept}

\section{Sequential Circuits}

\subsection{Flip-Flops}
\begin{concept}
\textbf{SR Flip-Flop}:
\begin{itemize}
    \item S=1, R=0: Set (Q=1)
    \item S=0, R=1: Reset (Q=0)
    \item S=0, R=0: Hold
    \item S=1, R=1: Invalid state
\end{itemize}

\textbf{D Flip-Flop}:
\begin{itemize}
    \item Q(next) = D
    \item Eliminates invalid state of SR flip-flop
    \item Used for data storage
\end{itemize}

\textbf{JK Flip-Flop}:
\begin{itemize}
    \item J=1, K=0: Set
    \item J=0, K=1: Reset
    \item J=0, K=0: Hold
    \item J=1, K=1: Toggle
\end{itemize}

\textbf{T Flip-Flop}:
\begin{itemize}
    \item T=0: Hold
    \item T=1: Toggle
    \item Used in counters
\end{itemize}
\end{concept}

\subsection{Registers and Counters}
\begin{concept}
\textbf{Registers}:
\begin{itemize}
    \item \textbf{Parallel Load}: All bits loaded simultaneously
    \item \textbf{Shift Register}: Data shifted left or right
    \item \textbf{SISO}: Serial In Serial Out
    \item \textbf{SIPO}: Serial In Parallel Out
    \item \textbf{PISO}: Parallel In Serial Out
    \item \textbf{PIPO}: Parallel In Parallel Out
\end{itemize}

\textbf{Counters}:
\begin{itemize}
    \item \textbf{Asynchronous (Ripple)}: Clock applied to first FF, output triggers next
    \item \textbf{Synchronous}: Common clock for all FFs
    \item \textbf{Up Counter}: Increments
    \item \textbf{Down Counter}: Decrements
    \item \textbf{Modulo-N Counter}: Counts from 0 to N-1
    \item \textbf{Ring Counter}: Circular shift register
    \item \textbf{Johnson Counter}: Twisted ring counter
\end{itemize}
\end{concept}

\section{Minimization Techniques}

\subsection{Karnaugh Maps (K-Maps)}
\begin{concept}
\textbf{K-Map Method}:
\begin{itemize}
    \item Visual method for minimizing Boolean functions
    \item Adjacent cells differ by one variable
    \item Group 1s in powers of 2 (1, 2, 4, 8, 16)
    \item Larger groups = simpler expressions
    \item Don't care conditions (X) can be treated as 0 or 1
\end{itemize}
\end{concept}

\subsection{Quine-McCluskey Method}
\begin{concept}
\textbf{Tabular Minimization}:
\begin{itemize}
    \item Systematic method for any number of variables
    \item Step 1: List all minterms
    \item Step 2: Group by number of 1s
    \item Step 3: Combine terms differing by one bit
    \item Step 4: Find prime implicants
    \item Step 5: Select essential prime implicants
\end{itemize}
\end{concept}

\chapter{Computer Organization and Architecture}

\section{Number Systems}

\subsection{Number Representations}
\begin{concept}
\textbf{Positional Number Systems}:
\begin{itemize}
    \item \textbf{Binary} (Base-2): Digits 0, 1
    \item \textbf{Octal} (Base-8): Digits 0-7
    \item \textbf{Decimal} (Base-10): Digits 0-9
    \item \textbf{Hexadecimal} (Base-16): Digits 0-9, A-F
\end{itemize}

\textbf{Conversions}:
\begin{itemize}
    \item Binary to Decimal: Multiply each bit by $2^i$
    \item Decimal to Binary: Repeated division by 2
    \item Binary to Hex: Group 4 bits
    \item Binary to Octal: Group 3 bits
\end{itemize}
\end{concept}

\subsection{Signed Number Representations}
\begin{concept}
\begin{itemize}
    \item \textbf{Sign-Magnitude}: MSB for sign, remaining for magnitude
    \item \textbf{1's Complement}: Flip all bits for negative
    \item \textbf{2's Complement}: Flip all bits and add 1
        \begin{itemize}
            \item Range: $-2^{n-1}$ to $2^{n-1} - 1$ for $n$ bits
            \item No separate representation for +0 and -0
            \item Addition/subtraction simplified
            \item Most widely used
        \end{itemize}
    \item \textbf{Excess/Bias Notation}: Add bias to actual value
\end{itemize}
\end{concept}

\subsection{Floating Point Representation}
\begin{concept}
\textbf{IEEE 754 Standard}:

\textbf{Single Precision (32 bits)}:
\begin{itemize}
    \item Sign: 1 bit
    \item Exponent: 8 bits (bias = 127)
    \item Mantissa: 23 bits
    \item Value: $(-1)^s \times 1.M \times 2^{E-127}$
\end{itemize}

\textbf{Double Precision (64 bits)}:
\begin{itemize}
    \item Sign: 1 bit
    \item Exponent: 11 bits (bias = 1023)
    \item Mantissa: 52 bits
    \item Value: $(-1)^s \times 1.M \times 2^{E-1023}$
\end{itemize}

\textbf{Special Values}:
\begin{itemize}
    \item All 0s exponent, 0 mantissa: 0
    \item All 1s exponent, 0 mantissa: Infinity
    \item All 1s exponent, non-zero mantissa: NaN
\end{itemize}
\end{concept}

\section{CPU Organization}

\subsection{Von Neumann Architecture}
\begin{concept}
\textbf{Components}:
\begin{itemize}
    \item \textbf{CPU}: Control Unit + ALU + Registers
    \item \textbf{Memory}: Stores both instructions and data
    \item \textbf{Input/Output}: Interfaces with external devices
    \item \textbf{Bus System}: Data bus, Address bus, Control bus
\end{itemize}

\textbf{Characteristics}:
\begin{itemize}
    \item Stored program concept
    \item Sequential instruction execution
    \item Single memory for instructions and data
    \item Von Neumann bottleneck: CPU-memory bandwidth limit
\end{itemize}
\end{concept}

\subsection{Harvard Architecture}
\begin{concept}
\begin{itemize}
    \item Separate memory for instructions and data
    \item Separate buses for instructions and data
    \item Can fetch instruction and data simultaneously
    \item Used in DSP processors and microcontrollers
\end{itemize}
\end{concept}

\subsection{Instruction Set Architecture (ISA)}
\begin{concept}
\textbf{CISC (Complex Instruction Set Computer)}:
\begin{itemize}
    \item Large instruction set
    \item Variable-length instructions
    \item Complex addressing modes
    \item More work per instruction
    \item Example: x86
\end{itemize}

\textbf{RISC (Reduced Instruction Set Computer)}:
\begin{itemize}
    \item Small instruction set
    \item Fixed-length instructions
    \item Simple addressing modes
    \item Load-store architecture
    \item More instructions per program
    \item Easier to pipeline
    \item Example: ARM, MIPS
\end{itemize}
\end{concept}

\section{Instruction Cycle}

\subsection{Fetch-Decode-Execute Cycle}
\begin{concept}
\textbf{Instruction Cycle Phases}:
\begin{enumerate}
    \item \textbf{Fetch}: Load instruction from memory using PC
        \begin{itemize}
            \item MAR $\leftarrow$ PC
            \item MBR $\leftarrow$ Memory[MAR]
            \item IR $\leftarrow$ MBR
            \item PC $\leftarrow$ PC + 1
        \end{itemize}
    \item \textbf{Decode}: Determine operation and operands
    \item \textbf{Execute}: Perform the operation
    \item \textbf{Store}: Write result back to memory/register
\end{enumerate}
\end{concept}

\subsection{Addressing Modes}
\begin{concept}
\begin{itemize}
    \item \textbf{Immediate}: Operand in instruction itself
    \item \textbf{Direct}: Address of operand in instruction
    \item \textbf{Indirect}: Instruction contains address of address
    \item \textbf{Register}: Operand in register
    \item \textbf{Register Indirect}: Register contains address of operand
    \item \textbf{Indexed}: Base address + index register
    \item \textbf{Base-Register}: Base register + offset in instruction
    \item \textbf{Relative}: PC + offset
    \item \textbf{Auto-increment/decrement}: Register value incremented/decremented
\end{itemize}
\end{concept}

\section{Pipelining}

\subsection{Pipeline Basics}
\begin{concept}
\textbf{Pipeline Stages (5-stage)}:
\begin{enumerate}
    \item \textbf{IF}: Instruction Fetch
    \item \textbf{ID}: Instruction Decode
    \item \textbf{EX}: Execute
    \item \textbf{MEM}: Memory Access
    \item \textbf{WB}: Write Back
\end{enumerate}

\textbf{Benefits}:
\begin{itemize}
    \item Increases throughput
    \item Multiple instructions in different stages simultaneously
    \item Ideal speedup = Number of stages
\end{itemize}
\end{concept}

\subsection{Pipeline Hazards}
\begin{concept}
\textbf{Structural Hazards}:
\begin{itemize}
    \item Hardware resource conflict
    \item Solution: Duplicate resources or stall
\end{itemize}

\textbf{Data Hazards}:
\begin{itemize}
    \item \textbf{RAW (Read After Write)}: True dependency
    \item \textbf{WAR (Write After Read)}: Anti-dependency
    \item \textbf{WAW (Write After Write)}: Output dependency
    \item Solutions: Forwarding, Stalling, Register renaming
\end{itemize}

\textbf{Control Hazards}:
\begin{itemize}
    \item Branch instructions change PC
    \item Solutions: Branch prediction, Delayed branching, Multiple streams
\end{itemize}
\end{concept}

\begin{formula}
\textbf{Pipeline Performance}:
\begin{itemize}
    \item \textbf{Speedup} = $\frac{\text{Time without pipeline}}{\text{Time with pipeline}}$
    \item \textbf{Throughput} = $\frac{\text{Number of instructions}}{\text{Time}}$
    \item \textbf{Efficiency} = $\frac{\text{Actual speedup}}{\text{Ideal speedup}}$
\end{itemize}
\end{formula}

\section{Memory Organization}

\subsection{Memory Hierarchy}
\begin{concept}
\textbf{Hierarchy Levels} (Fast to Slow):
\begin{enumerate}
    \item Registers
    \item Cache (L1, L2, L3)
    \item Main Memory (RAM)
    \item Secondary Storage (HDD, SSD)
\end{enumerate}

\textbf{Principles}:
\begin{itemize}
    \item \textbf{Temporal Locality}: Recently accessed items likely to be accessed again
    \item \textbf{Spatial Locality}: Items near recently accessed items likely to be accessed
\end{itemize}
\end{concept}

\subsection{Cache Memory}
\begin{concept}
\textbf{Cache Mapping Techniques}:

\textbf{Direct Mapped}:
\begin{itemize}
    \item Each block maps to exactly one cache line
    \item Cache line = (Block address) mod (Number of cache lines)
    \item Fast but high conflict misses
\end{itemize}

\textbf{Fully Associative}:
\begin{itemize}
    \item Block can go in any cache line
    \item Flexible but expensive (requires parallel search)
\end{itemize}

\textbf{Set Associative}:
\begin{itemize}
    \item N-way set associative: N blocks per set
    \item Block maps to a set, can go in any line within set
    \item Balance between direct and fully associative
\end{itemize}
\end{concept}

\subsection{Cache Replacement Policies}
\begin{concept}
\begin{itemize}
    \item \textbf{LRU (Least Recently Used)}: Replace block unused for longest time
    \item \textbf{FIFO (First In First Out)}: Replace oldest block
    \item \textbf{LFU (Least Frequently Used)}: Replace block with lowest access count
    \item \textbf{Random}: Replace random block
\end{itemize}
\end{concept}

\subsection{Cache Write Policies}
\begin{concept}
\textbf{Write Hit}:
\begin{itemize}
    \item \textbf{Write Through}: Write to both cache and memory
    \item \textbf{Write Back}: Write only to cache, update memory on replacement
\end{itemize}

\textbf{Write Miss}:
\begin{itemize}
    \item \textbf{Write Allocate}: Load block into cache, then write
    \item \textbf{No Write Allocate}: Write directly to memory
\end{itemize}
\end{concept}

\begin{formula}
\textbf{Cache Performance}:
\begin{itemize}
    \item \textbf{Hit Rate} = $\frac{\text{Number of hits}}{\text{Total accesses}}$
    \item \textbf{Miss Rate} = $1 - \text{Hit Rate}$
    \item \textbf{Average Access Time} = Hit Time + (Miss Rate $\times$ Miss Penalty)
\end{itemize}
\end{formula}

\subsection{Virtual Memory}
\begin{concept}
\textbf{Virtual Memory Concepts}:
\begin{itemize}
    \item Abstraction of physical memory
    \item Provides large address space
    \item Allows multiprogramming
    \item Memory protection between processes
\end{itemize}

\textbf{Paging}:
\begin{itemize}
    \item Virtual and physical memory divided into fixed-size pages
    \item \textbf{Page Table}: Maps virtual to physical pages
    \item \textbf{Page Fault}: Requested page not in memory
    \item \textbf{TLB (Translation Lookaside Buffer)}: Cache for page table entries
\end{itemize}

\textbf{Segmentation}:
\begin{itemize}
    \item Memory divided into variable-size segments
    \item Logical division (code, data, stack)
    \item \textbf{Segment Table}: Base address and limit for each segment
\end{itemize}
\end{concept}

\subsection{Page Replacement Algorithms}
\begin{concept}
\begin{itemize}
    \item \textbf{FIFO}: Replace oldest page
    \item \textbf{Optimal}: Replace page not used for longest time in future (theoretical)
    \item \textbf{LRU}: Replace least recently used page
    \item \textbf{Clock/Second Chance}: FIFO with reference bit
    \item \textbf{LFU}: Replace least frequently used page
\end{itemize}

\textbf{Belady's Anomaly}: More frames can lead to more page faults (FIFO)
\end{concept}

\section{Input/Output Organization}

\subsection{I/O Techniques}
\begin{concept}
\textbf{Programmed I/O}:
\begin{itemize}
    \item CPU polls device status
    \item CPU waits for I/O completion
    \item Simple but wastes CPU cycles
\end{itemize}

\textbf{Interrupt-Driven I/O}:
\begin{itemize}
    \item Device signals CPU when ready
    \item CPU can do other work while waiting
    \item Context switching overhead
\end{itemize}

\textbf{Direct Memory Access (DMA)}:
\begin{itemize}
    \item DMA controller transfers data directly to/from memory
    \item CPU only involved at start and end
    \item Most efficient for large transfers
    \item Cycle stealing: DMA temporarily takes bus control
\end{itemize}
\end{concept}

\subsection{I/O Performance}
\begin{formula}
\textbf{I/O Metrics}:
\begin{itemize}
    \item \textbf{Throughput}: Data transferred per unit time
    \item \textbf{Latency}: Time for single I/O operation
    \item \textbf{Bandwidth}: Maximum data transfer rate
\end{itemize}
\end{formula}

\chapter{Programming and Data Structures}

\section{Programming Basics}

\subsection{Algorithm Analysis}
\begin{concept}
\textbf{Time Complexity}:
\begin{itemize}
    \item Measures growth rate of execution time
    \item Expressed using Big-O notation
    \item Worst case, Average case, Best case
\end{itemize}

\textbf{Space Complexity}:
\begin{itemize}
    \item Memory required by algorithm
    \item Includes auxiliary space and input space
\end{itemize}
\end{concept}

\subsection{Asymptotic Notations}
\begin{concept}
\begin{itemize}
    \item \textbf{Big-O} ($O$): Upper bound (worst case)
        \begin{itemize}
            \item $f(n) = O(g(n))$ if $\exists c, n_0$ such that $f(n) \leq c \cdot g(n)$ for all $n \geq n_0$
        \end{itemize}
    \item \textbf{Big-Omega} ($\Omega$): Lower bound (best case)
        \begin{itemize}
            \item $f(n) = \Omega(g(n))$ if $\exists c, n_0$ such that $f(n) \geq c \cdot g(n)$ for all $n \geq n_0$
        \end{itemize}
    \item \textbf{Big-Theta} ($\Theta$): Tight bound (average case)
        \begin{itemize}
            \item $f(n) = \Theta(g(n))$ if $f(n) = O(g(n))$ and $f(n) = \Omega(g(n))$
        \end{itemize}
    \item \textbf{Little-o} ($o$): Strict upper bound
    \item \textbf{Little-omega} ($\omega$): Strict lower bound
\end{itemize}
\end{concept}

\subsection{Common Time Complexities}
\begin{concept}
\begin{itemize}
    \item $O(1)$: Constant - Array access, hash table operations
    \item $O(\log n)$: Logarithmic - Binary search, balanced BST operations
    \item $O(n)$: Linear - Linear search, array traversal
    \item $O(n \log n)$: Linearithmic - Merge sort, heap sort, quick sort (average)
    \item $O(n^2)$: Quadratic - Bubble sort, selection sort, insertion sort
    \item $O(n^3)$: Cubic - Matrix multiplication (naive)
    \item $O(2^n)$: Exponential - Recursive Fibonacci, subset generation
    \item $O(n!)$: Factorial - Permutation generation
\end{itemize}
\end{concept}

\section{Arrays}

\subsection{Array Operations}
\begin{concept}
\textbf{Operations and Complexity}:
\begin{itemize}
    \item Access: $O(1)$
    \item Search: $O(n)$ unsorted, $O(\log n)$ sorted
    \item Insertion: $O(n)$ (shifting required)
    \item Deletion: $O(n)$ (shifting required)
    \item Update: $O(1)$
\end{itemize}

\textbf{Advantages}:
\begin{itemize}
    \item Random access
    \item Cache friendly
    \item Simple implementation
\end{itemize}

\textbf{Disadvantages}:
\begin{itemize}
    \item Fixed size
    \item Expensive insertion/deletion
\end{itemize}
\end{concept}

\section{Linked Lists}

\subsection{Types of Linked Lists}
\begin{concept}
\textbf{Singly Linked List}:
\begin{itemize}
    \item Each node has data and next pointer
    \item One-way traversal
\end{itemize}

\textbf{Doubly Linked List}:
\begin{itemize}
    \item Each node has data, next, and prev pointers
    \item Bi-directional traversal
    \item Easier deletion but more space
\end{itemize}

\textbf{Circular Linked List}:
\begin{itemize}
    \item Last node points to first node
    \item Can be singly or doubly linked
\end{itemize}
\end{concept}

\subsection{Linked List Operations}
\begin{concept}
\textbf{Time Complexity}:
\begin{itemize}
    \item Access: $O(n)$
    \item Search: $O(n)$
    \item Insertion at beginning: $O(1)$
    \item Insertion at end: $O(n)$ for singly, $O(1)$ with tail pointer
    \item Deletion (node given): $O(1)$ for doubly, $O(n)$ for singly
\end{itemize}

\textbf{Advantages}:
\begin{itemize}
    \item Dynamic size
    \item Efficient insertion/deletion
\end{itemize}

\textbf{Disadvantages}:
\begin{itemize}
    \item No random access
    \item Extra memory for pointers
    \item Not cache friendly
\end{itemize}
\end{concept}

\section{Stacks}

\subsection{Stack Operations}
\begin{concept}
\textbf{LIFO (Last In First Out)}:
\begin{itemize}
    \item \textbf{Push}: Add element to top - $O(1)$
    \item \textbf{Pop}: Remove element from top - $O(1)$
    \item \textbf{Peek/Top}: View top element - $O(1)$
    \item \textbf{isEmpty}: Check if empty - $O(1)$
\end{itemize}

\textbf{Applications}:
\begin{itemize}
    \item Function call stack
    \item Expression evaluation
    \item Backtracking algorithms
    \item Undo mechanisms
    \item Parenthesis matching
\end{itemize}
\end{concept}

\subsection{Expression Evaluation}
\begin{concept}
\textbf{Infix, Prefix, Postfix}:
\begin{itemize}
    \item \textbf{Infix}: $A + B$ (operator between operands)
    \item \textbf{Prefix/Polish}: $+ A B$ (operator before operands)
    \item \textbf{Postfix/Reverse Polish}: $A B +$ (operator after operands)
\end{itemize}

\textbf{Conversion}: Use stack to convert between notations

\textbf{Evaluation}: Postfix and prefix are easier to evaluate using stack
\end{concept}

\section{Queues}

\subsection{Queue Operations}
\begin{concept}
\textbf{FIFO (First In First Out)}:
\begin{itemize}
    \item \textbf{Enqueue}: Add element to rear - $O(1)$
    \item \textbf{Dequeue}: Remove element from front - $O(1)$
    \item \textbf{Front}: View front element - $O(1)$
    \item \textbf{isEmpty}: Check if empty - $O(1)$
\end{itemize}

\textbf{Types}:
\begin{itemize}
    \item \textbf{Simple Queue}: Linear queue
    \item \textbf{Circular Queue}: Last position connects to first
    \item \textbf{Priority Queue}: Elements have priorities
    \item \textbf{Deque}: Double-ended queue (insertion/deletion at both ends)
\end{itemize}
\end{concept}

\section{Trees}

\subsection{Binary Trees}
\begin{concept}
\textbf{Binary Tree}: Each node has at most 2 children

\textbf{Types}:
\begin{itemize}
    \item \textbf{Full Binary Tree}: Every node has 0 or 2 children
    \item \textbf{Complete Binary Tree}: All levels filled except possibly last, filled left to right
    \item \textbf{Perfect Binary Tree}: All internal nodes have 2 children, all leaves at same level
    \item \textbf{Balanced Binary Tree}: Height difference between left and right subtrees $\leq 1$
\end{itemize}
\end{concept}

\begin{formula}
\textbf{Binary Tree Properties}:
\begin{itemize}
    \item Maximum nodes at level $l$: $2^l$
    \item Maximum nodes in tree of height $h$: $2^{h+1} - 1$
    \item Minimum height for $n$ nodes: $\lceil \log_2(n+1) \rceil - 1$
    \item For $n$ nodes: $n = n_0 + n_1 + n_2$ where $n_i$ is nodes with $i$ children
    \item For binary tree: $n_0 = n_2 + 1$ (leaf nodes = internal nodes with 2 children + 1)
\end{itemize}
\end{formula}

\subsection{Tree Traversals}
\begin{concept}
\textbf{Depth First Traversals}:
\begin{itemize}
    \item \textbf{Inorder} (Left-Root-Right): For BST, gives sorted order
    \item \textbf{Preorder} (Root-Left-Right): Used to copy tree
    \item \textbf{Postorder} (Left-Right-Root): Used to delete tree
\end{itemize}

\textbf{Breadth First Traversal}:
\begin{itemize}
    \item \textbf{Level Order}: Visit nodes level by level (use queue)
\end{itemize}

All traversals: Time $O(n)$, Space $O(h)$ where $h$ is height
\end{concept}

\subsection{Binary Search Tree (BST)}
\begin{concept}
\textbf{BST Property}: For each node:
\begin{itemize}
    \item All nodes in left subtree $<$ node value
    \item All nodes in right subtree $>$ node value
\end{itemize}

\textbf{Operations}:
\begin{itemize}
    \item Search: $O(h)$ where $h$ is height
    \item Insertion: $O(h)$
    \item Deletion: $O(h)$
    \item Inorder traversal gives sorted sequence
\end{itemize}

\textbf{Worst Case}: Skewed tree, $h = n$, all operations $O(n)$
\textbf{Best Case}: Balanced tree, $h = \log n$, all operations $O(\log n)$
\end{concept}

\subsection{AVL Trees}
\begin{concept}
\textbf{Self-Balancing BST}:
\begin{itemize}
    \item For every node: $|height(left) - height(right)| \leq 1$
    \item \textbf{Balance Factor}: height(left) - height(right) $\in \{-1, 0, 1\}$
\end{itemize}

\textbf{Rotations}:
\begin{itemize}
    \item \textbf{LL Rotation}: Right rotation
    \item \textbf{RR Rotation}: Left rotation
    \item \textbf{LR Rotation}: Left-Right rotation
    \item \textbf{RL Rotation}: Right-Left rotation
\end{itemize}

\textbf{Complexity}: All operations $O(\log n)$
\end{concept}

\subsection{Red-Black Trees}
\begin{concept}
\textbf{Properties}:
\begin{enumerate}
    \item Every node is either red or black
    \item Root is black
    \item All leaves (NIL) are black
    \item If a node is red, both children are black
    \item All paths from node to descendant leaves have same number of black nodes
\end{enumerate}

\textbf{Complexity}: All operations $O(\log n)$

\textbf{Comparison with AVL}:
\begin{itemize}
    \item AVL more rigidly balanced, faster lookups
    \item Red-Black trees faster insertion/deletion
\end{itemize}
\end{concept}

\subsection{B-Trees}
\begin{concept}
\textbf{B-Tree of order $m$}:
\begin{itemize}
    \item Each node has at most $m$ children
    \item Non-leaf nodes (except root) have at least $\lceil m/2 \rceil$ children
    \item Root has at least 2 children (unless leaf)
    \item All leaves at same level
    \item Node with $k$ children has $k-1$ keys
\end{itemize}

\textbf{Applications}: Databases, file systems

\textbf{Complexity}: Search, Insert, Delete all $O(\log n)$
\end{concept}

\section{Heaps}

\subsection{Binary Heap}
\begin{concept}
\textbf{Heap Property}:
\begin{itemize}
    \item \textbf{Max Heap}: Parent $\geq$ children
    \item \textbf{Min Heap}: Parent $\leq$ children
\end{itemize}

\textbf{Complete Binary Tree Structure}: Can be stored in array
\begin{itemize}
    \item Parent of node $i$: $\lfloor (i-1)/2 \rfloor$
    \item Left child of node $i$: $2i + 1$
    \item Right child of node $i$: $2i + 2$
\end{itemize}

\textbf{Operations}:
\begin{itemize}
    \item Get Max/Min: $O(1)$
    \item Insert: $O(\log n)$ (heapify up)
    \item Delete Max/Min: $O(\log n)$ (heapify down)
    \item Build Heap: $O(n)$
\end{itemize}

\textbf{Applications}: Priority Queue, Heap Sort
\end{concept}

\section{Hashing}

\subsection{Hash Tables}
\begin{concept}
\textbf{Hash Function}: Maps key to index in array

\textbf{Properties of Good Hash Function}:
\begin{itemize}
    \item Uniform distribution
    \item Fast computation
    \item Minimize collisions
\end{itemize}

\textbf{Common Hash Functions}:
\begin{itemize}
    \item Division Method: $h(k) = k \mod m$
    \item Multiplication Method: $h(k) = \lfloor m(kA \mod 1) \rfloor$
    \item Universal Hashing
\end{itemize}
\end{concept}

\subsection{Collision Resolution}
\begin{concept}
\textbf{Open Addressing}:
\begin{itemize}
    \item \textbf{Linear Probing}: $h(k,i) = (h'(k) + i) \mod m$
        \begin{itemize}
            \item Simple but causes clustering
        \end{itemize}
    \item \textbf{Quadratic Probing}: $h(k,i) = (h'(k) + c_1 i + c_2 i^2) \mod m$
        \begin{itemize}
            \item Reduces primary clustering
        \end{itemize}
    \item \textbf{Double Hashing}: $h(k,i) = (h_1(k) + i \cdot h_2(k)) \mod m$
        \begin{itemize}
            \item Best distribution
        \end{itemize}
\end{itemize}

\textbf{Chaining}:
\begin{itemize}
    \item Each bucket stores linked list of colliding elements
    \item Simple to implement
    \item Load factor can exceed 1
\end{itemize}
\end{concept}

\begin{formula}
\textbf{Load Factor}: $\alpha = \frac{n}{m}$ where $n$ is number of elements, $m$ is table size

\textbf{Expected Time Complexity}:
\begin{itemize}
    \item Chaining: $O(1 + \alpha)$
    \item Open Addressing: $O(1/(1-\alpha))$ for unsuccessful search
\end{itemize}
\end{formula}

\section{Graphs}

\subsection{Graph Representations}
\begin{concept}
\textbf{Adjacency Matrix}:
\begin{itemize}
    \item 2D array $A[i][j] = 1$ if edge from $i$ to $j$
    \item Space: $O(V^2)$
    \item Check edge: $O(1)$
    \item Find neighbors: $O(V)$
    \item Good for dense graphs
\end{itemize}

\textbf{Adjacency List}:
\begin{itemize}
    \item Array of lists, each list stores neighbors
    \item Space: $O(V + E)$
    \item Check edge: $O(degree)$
    \item Find neighbors: $O(degree)$
    \item Good for sparse graphs
\end{itemize}
\end{concept}

\subsection{Graph Traversal}
\begin{concept}
\textbf{Breadth First Search (BFS)}:
\begin{itemize}
    \item Uses queue
    \item Explores level by level
    \item Time: $O(V + E)$
    \item Space: $O(V)$
    \item Finds shortest path in unweighted graph
\end{itemize}

\textbf{Depth First Search (DFS)}:
\begin{itemize}
    \item Uses stack (or recursion)
    \item Explores as far as possible before backtracking
    \item Time: $O(V + E)$
    \item Space: $O(V)$
    \item Used in topological sorting, cycle detection, strongly connected components
\end{itemize}
\end{concept}

\subsection{Topological Sorting}
\begin{concept}
\textbf{Topological Sort}:
\begin{itemize}
    \item Linear ordering of vertices in DAG
    \item For edge $u \to v$, $u$ comes before $v$
    \item Exists only for DAGs (Directed Acyclic Graphs)
\end{itemize}

\textbf{Algorithms}:
\begin{itemize}
    \item DFS-based: $O(V + E)$
    \item Kahn's Algorithm (BFS): $O(V + E)$
\end{itemize}
\end{concept}

\chapter{Algorithms}

\section{Sorting Algorithms}

\subsection{Comparison-Based Sorting}
\begin{concept}
\textbf{Bubble Sort}:
\begin{itemize}
    \item Repeatedly swap adjacent elements if in wrong order
    \item Time: $O(n^2)$ worst/average, $O(n)$ best
    \item Space: $O(1)$
    \item Stable
\end{itemize}

\textbf{Selection Sort}:
\begin{itemize}
    \item Find minimum and place at beginning
    \item Time: $O(n^2)$ all cases
    \item Space: $O(1)$
    \item Not stable (can be made stable)
\end{itemize}

\textbf{Insertion Sort}:
\begin{itemize}
    \item Insert each element in its correct position in sorted part
    \item Time: $O(n^2)$ worst/average, $O(n)$ best
    \item Space: $O(1)$
    \item Stable
    \item Good for nearly sorted data
\end{itemize}
\end{concept}

\begin{concept}
\textbf{Merge Sort}:
\begin{itemize}
    \item Divide and conquer
    \item Recursively divide, then merge sorted halves
    \item Time: $O(n \log n)$ all cases
    \item Space: $O(n)$
    \item Stable
    \item External sorting
\end{itemize}

\textbf{Quick Sort}:
\begin{itemize}
    \item Choose pivot, partition around pivot
    \item Time: $O(n^2)$ worst, $O(n \log n)$ average/best
    \item Space: $O(\log n)$ stack space
    \item Not stable
    \item In-place
    \item Worst case when already sorted (use random pivot)
\end{itemize}

\textbf{Heap Sort}:
\begin{itemize}
    \item Build max heap, repeatedly extract maximum
    \item Time: $O(n \log n)$ all cases
    \item Space: $O(1)$
    \item Not stable
    \item In-place
\end{itemize}
\end{concept}

\subsection{Non-Comparison Sorting}
\begin{concept}
\textbf{Counting Sort}:
\begin{itemize}
    \item Count occurrences of each value
    \item Time: $O(n + k)$ where $k$ is range
    \item Space: $O(k)$
    \item Stable
    \item Works when range is small
\end{itemize}

\textbf{Radix Sort}:
\begin{itemize}
    \item Sort digit by digit using stable sort
    \item Time: $O(d(n + k))$ where $d$ is digits, $k$ is base
    \item Space: $O(n + k)$
    \item Stable
\end{itemize}

\textbf{Bucket Sort}:
\begin{itemize}
    \item Distribute elements into buckets, sort each bucket
    \item Time: $O(n)$ average, $O(n^2)$ worst
    \item Space: $O(n + k)$
    \item Works well for uniformly distributed data
\end{itemize}
\end{concept}

\begin{theorem}
\textbf{Lower Bound for Comparison Sorting}: Any comparison-based sorting algorithm requires $\Omega(n \log n)$ comparisons in the worst case.
\end{theorem}

\section{Searching Algorithms}

\subsection{Linear and Binary Search}
\begin{concept}
\textbf{Linear Search}:
\begin{itemize}
    \item Check each element sequentially
    \item Time: $O(n)$
    \item Works on unsorted arrays
\end{itemize}

\textbf{Binary Search}:
\begin{itemize}
    \item Divide search space in half
    \item Time: $O(\log n)$
    \item Requires sorted array
    \item Can be implemented iteratively or recursively
\end{itemize}
\end{concept}

\section{Divide and Conquer}

\subsection{Paradigm}
\begin{concept}
\textbf{Strategy}:
\begin{enumerate}
    \item \textbf{Divide}: Break problem into smaller subproblems
    \item \textbf{Conquer}: Solve subproblems recursively
    \item \textbf{Combine}: Merge solutions
\end{enumerate}

\textbf{Examples}:
\begin{itemize}
    \item Merge Sort
    \item Quick Sort
    \item Binary Search
    \item Strassen's Matrix Multiplication
\end{itemize}
\end{concept}

\subsection{Master Theorem}
\begin{theorem}
For recurrence $T(n) = aT(n/b) + f(n)$ where $a \geq 1$, $b > 1$:

Let $c_{crit} = \log_b a$

\begin{enumerate}
    \item If $f(n) = O(n^c)$ where $c < c_{crit}$: $T(n) = \Theta(n^{c_{crit}})$
    \item If $f(n) = \Theta(n^{c_{crit}} \log^k n)$: $T(n) = \Theta(n^{c_{crit}} \log^{k+1} n)$
    \item If $f(n) = \Omega(n^c)$ where $c > c_{crit}$ and $af(n/b) \leq kf(n)$ for some $k < 1$: $T(n) = \Theta(f(n))$
\end{enumerate}
\end{theorem}

\section{Greedy Algorithms}

\subsection{Greedy Strategy}
\begin{concept}
\textbf{Greedy Approach}:
\begin{itemize}
    \item Make locally optimal choice at each step
    \item Hope for globally optimal solution
    \item Does not always work
    \item When it works, very efficient
\end{itemize}

\textbf{Greedy Choice Property}: Locally optimal choice leads to globally optimal solution

\textbf{Optimal Substructure}: Optimal solution contains optimal solutions to subproblems
\end{concept}

\subsection{Classic Greedy Algorithms}
\begin{concept}
\textbf{Activity Selection}:
\begin{itemize}
    \item Select maximum number of non-overlapping activities
    \item Sort by finish time, greedily select
    \item Time: $O(n \log n)$
\end{itemize}

\textbf{Fractional Knapsack}:
\begin{itemize}
    \item Maximize value with weight constraint
    \item Sort by value/weight ratio
    \item Time: $O(n \log n)$
\end{itemize}

\textbf{Huffman Coding}:
\begin{itemize}
    \item Optimal prefix-free encoding
    \item Build tree bottom-up using min heap
    \item Time: $O(n \log n)$
\end{itemize}
\end{concept}

\subsection{Minimum Spanning Tree}
\begin{concept}
\textbf{Prim's Algorithm}:
\begin{itemize}
    \item Start from arbitrary vertex
    \item Add minimum weight edge connecting tree to non-tree vertex
    \item Time: $O(E \log V)$ with min heap
    \item Works on connected graphs
\end{itemize}

\textbf{Kruskal's Algorithm}:
\begin{itemize}
    \item Sort edges by weight
    \item Add edge if doesn't create cycle
    \item Use union-find for cycle detection
    \item Time: $O(E \log E)$ or $O(E \log V)$
\end{itemize}
\end{concept}

\subsection{Shortest Path}
\begin{concept}
\textbf{Dijkstra's Algorithm}:
\begin{itemize}
    \item Single source shortest path
    \item Non-negative edge weights
    \item Use priority queue
    \item Time: $O((V + E) \log V)$ with binary heap
    \item Greedy approach
\end{itemize}

\textbf{Bellman-Ford Algorithm}:
\begin{itemize}
    \item Single source shortest path
    \item Handles negative weights
    \item Detects negative cycles
    \item Time: $O(VE)$
    \item Dynamic programming approach
\end{itemize}

\textbf{Floyd-Warshall Algorithm}:
\begin{itemize}
    \item All pairs shortest path
    \item Time: $O(V^3)$
    \item Space: $O(V^2)$
    \item Dynamic programming
\end{itemize}
\end{concept}

\section{Dynamic Programming}

\subsection{DP Paradigm}
\begin{concept}
\textbf{Dynamic Programming}:
\begin{itemize}
    \item Break problem into overlapping subproblems
    \item Store solutions to avoid recomputation
    \item Bottom-up (tabulation) or Top-down (memoization)
\end{itemize}

\textbf{Requirements}:
\begin{itemize}
    \item \textbf{Optimal Substructure}: Optimal solution contains optimal substructure solutions
    \item \textbf{Overlapping Subproblems}: Same subproblems solved multiple times
\end{itemize}

\textbf{Steps}:
\begin{enumerate}
    \item Define state
    \item Find recurrence relation
    \item Identify base cases
    \item Compute bottom-up or top-down
\end{enumerate}
\end{concept}

\subsection{Classic DP Problems}
\begin{concept}
\textbf{Fibonacci Numbers}:
\begin{itemize}
    \item $F(n) = F(n-1) + F(n-2)$
    \item Time: $O(n)$, Space: $O(n)$ or $O(1)$
\end{itemize}

\textbf{0/1 Knapsack}:
\begin{itemize}
    \item Maximize value without exceeding weight
    \item Can't take fractions
    \item Time: $O(nW)$, Space: $O(nW)$
\end{itemize}

\textbf{Longest Common Subsequence (LCS)}:
\begin{itemize}
    \item Find longest subsequence common to two sequences
    \item Time: $O(mn)$, Space: $O(mn)$
\end{itemize}

\textbf{Longest Increasing Subsequence (LIS)}:
\begin{itemize}
    \item Find longest strictly increasing subsequence
    \item Time: $O(n^2)$ DP, $O(n \log n)$ with binary search
\end{itemize}

\textbf{Edit Distance}:
\begin{itemize}
    \item Minimum operations to convert one string to another
    \item Operations: Insert, Delete, Replace
    \item Time: $O(mn)$, Space: $O(mn)$
\end{itemize}

\textbf{Matrix Chain Multiplication}:
\begin{itemize}
    \item Optimal parenthesization
    \item Time: $O(n^3)$, Space: $O(n^2)$
\end{itemize}
\end{concept}

\section{Backtracking}

\subsection{Backtracking Paradigm}
\begin{concept}
\textbf{Backtracking}:
\begin{itemize}
    \item Try all possibilities recursively
    \item Abandon path if it can't lead to solution
    \item Implicit tree traversal
\end{itemize}

\textbf{When to Use}:
\begin{itemize}
    \item All solutions needed
    \item Optimization not primary goal
    \item Constraints can prune search space
\end{itemize}

\textbf{Examples}:
\begin{itemize}
    \item N-Queens problem
    \item Sudoku solver
    \item Graph coloring
    \item Hamiltonian path
    \item Subset sum
\end{itemize}
\end{concept}

\section{String Algorithms}

\subsection{Pattern Matching}
\begin{concept}
\textbf{Naive String Matching}:
\begin{itemize}
    \item Check pattern at each position
    \item Time: $O(nm)$ where $n$ is text length, $m$ is pattern length
\end{itemize}

\textbf{KMP (Knuth-Morris-Pratt)}:
\begin{itemize}
    \item Preprocess pattern to find longest prefix suffix
    \item Time: $O(n + m)$
    \item Space: $O(m)$
\end{itemize}

\textbf{Rabin-Karp}:
\begin{itemize}
    \item Rolling hash for pattern matching
    \item Time: $O(n + m)$ average, $O(nm)$ worst
    \item Good for multiple pattern matching
\end{itemize}
\end{concept}

\chapter{Theory of Computation}

\section{Finite Automata}

\subsection{Deterministic Finite Automaton (DFA)}
\begin{concept}
\textbf{DFA} = $(Q, \Sigma, \delta, q_0, F)$
\begin{itemize}
    \item $Q$: Finite set of states
    \item $\Sigma$: Input alphabet
    \item $\delta: Q \times \Sigma \to Q$: Transition function
    \item $q_0 \in Q$: Start state
    \item $F \subseteq Q$: Set of accepting states
\end{itemize}

\textbf{Properties}:
\begin{itemize}
    \item Exactly one transition for each state-symbol pair
    \item No $\epsilon$-transitions
    \item Deterministic behavior
\end{itemize}
\end{concept}

\subsection{Non-deterministic Finite Automaton (NFA)}
\begin{concept}
\textbf{NFA} = $(Q, \Sigma, \delta, q_0, F)$
\begin{itemize}
    \item $\delta: Q \times (\Sigma \cup \{\epsilon\}) \to 2^Q$: Transition function to set of states
\end{itemize}

\textbf{Properties}:
\begin{itemize}
    \item Multiple transitions possible for same state-symbol
    \item $\epsilon$-transitions allowed
    \item Non-deterministic behavior
\end{itemize}

\textbf{Relationship}:
\begin{itemize}
    \item Every DFA is an NFA
    \item Every NFA can be converted to equivalent DFA (subset construction)
    \item DFA and NFA have same expressive power (recognize regular languages)
\end{itemize}
\end{concept}

\section{Regular Languages}

\subsection{Regular Expressions}
\begin{concept}
\textbf{Regular Expression Operators}:
\begin{itemize}
    \item \textbf{Union}: $r_1 + r_2$ or $r_1 | r_2$
    \item \textbf{Concatenation}: $r_1 \cdot r_2$ or $r_1 r_2$
    \item \textbf{Kleene Star}: $r^*$ (zero or more repetitions)
    \item \textbf{Plus}: $r^+$ (one or more repetitions)
\end{itemize}

\textbf{Equivalences}:
\begin{itemize}
    \item Regular Expression $\Leftrightarrow$ NFA $\Leftrightarrow$ DFA
    \item All represent regular languages
\end{itemize}
\end{concept}

\subsection{Properties of Regular Languages}
\begin{concept}
\textbf{Closure Properties} (Regular languages closed under):
\begin{itemize}
    \item Union
    \item Concatenation
    \item Kleene Star
    \item Intersection
    \item Complement
    \item Difference
    \item Reversal
\end{itemize}

\textbf{Not Closed Under}:
\begin{itemize}
    \item None (regular languages have very strong closure properties)
\end{itemize}
\end{concept}

\subsection{Pumping Lemma for Regular Languages}
\begin{theorem}
\textbf{Pumping Lemma}: If $L$ is regular, then $\exists p \geq 1$ (pumping length) such that $\forall$ string $s \in L$ with $|s| \geq p$, $s$ can be divided into $s = xyz$ where:
\begin{enumerate}
    \item $|y| > 0$
    \item $|xy| \leq p$
    \item $\forall i \geq 0: xy^iz \in L$
\end{enumerate}

\textbf{Use}: To prove a language is NOT regular (proof by contradiction)
\end{theorem}

\section{Context-Free Languages}

\subsection{Context-Free Grammar (CFG)}
\begin{concept}
\textbf{CFG} = $(V, \Sigma, R, S)$
\begin{itemize}
    \item $V$: Set of variables (non-terminals)
    \item $\Sigma$: Set of terminals
    \item $R$: Set of production rules ($A \to \alpha$ where $A \in V$, $\alpha \in (V \cup \Sigma)^*$)
    \item $S$: Start symbol
\end{itemize}

\textbf{Derivations}:
\begin{itemize}
    \item \textbf{Leftmost}: Always expand leftmost variable first
    \item \textbf{Rightmost}: Always expand rightmost variable first
\end{itemize}

\textbf{Parse Tree}: Tree representation of derivation
\end{concept}

\subsection{Ambiguity}
\begin{concept}
\textbf{Ambiguous Grammar}: Grammar with more than one parse tree (or leftmost/rightmost derivation) for some string

\textbf{Inherently Ambiguous Language}: Every CFG for the language is ambiguous

\textbf{Resolving Ambiguity}:
\begin{itemize}
    \item Operator precedence
    \item Associativity rules
    \item Grammar rewriting
\end{itemize}
\end{concept}

\subsection{Normal Forms}
\begin{concept}
\textbf{Chomsky Normal Form (CNF)}:
\begin{itemize}
    \item Productions: $A \to BC$ or $A \to a$
    \item Where $A, B, C$ are variables, $a$ is terminal
    \item Every CFG (without $\epsilon$) can be converted to CNF
\end{itemize}

\textbf{Greibach Normal Form (GNF)}:
\begin{itemize}
    \item Productions: $A \to a\alpha$
    \item Where $a$ is terminal, $\alpha$ is string of variables
\end{itemize}
\end{concept}

\subsection{Pushdown Automaton (PDA)}
\begin{concept}
\textbf{PDA}: Finite automaton + stack

\textbf{PDA} = $(Q, \Sigma, \Gamma, \delta, q_0, Z_0, F)$
\begin{itemize}
    \item $Q$: States
    \item $\Sigma$: Input alphabet
    \item $\Gamma$: Stack alphabet
    \item $\delta: Q \times (\Sigma \cup \{\epsilon\}) \times \Gamma \to 2^{Q \times \Gamma^*}$
    \item $q_0$: Start state
    \item $Z_0$: Initial stack symbol
    \item $F$: Accepting states
\end{itemize}

\textbf{Equivalence}:
\begin{itemize}
    \item PDA $\Leftrightarrow$ CFG
    \item Both recognize context-free languages
\end{itemize}
\end{concept}

\subsection{Pumping Lemma for CFLs}
\begin{theorem}
\textbf{Pumping Lemma for CFLs}: If $L$ is CFL, then $\exists p \geq 1$ such that $\forall$ string $s \in L$ with $|s| \geq p$, $s$ can be divided into $s = uvwxy$ where:
\begin{enumerate}
    \item $|vx| > 0$
    \item $|vwx| \leq p$
    \item $\forall i \geq 0: uv^iwx^iy \in L$
\end{enumerate}

\textbf{Use}: To prove a language is NOT context-free
\end{theorem}

\section{Turing Machines}

\subsection{Turing Machine Definition}
\begin{concept}
\textbf{Turing Machine (TM)} = $(Q, \Sigma, \Gamma, \delta, q_0, q_{accept}, q_{reject})$
\begin{itemize}
    \item $Q$: Finite set of states
    \item $\Sigma$: Input alphabet (not including blank)
    \item $\Gamma$: Tape alphabet ($\Sigma \subseteq \Gamma$)
    \item $\delta: Q \times \Gamma \to Q \times \Gamma \times \{L, R\}$
    \item $q_0$: Start state
    \item $q_{accept}$: Accept state
    \item $q_{reject}$: Reject state
\end{itemize}

\textbf{Properties}:
\begin{itemize}
    \item Infinite tape (both directions)
    \item Can read and write on tape
    \item Can move left or right
\end{itemize}
\end{concept}

\subsection{Variants of Turing Machines}
\begin{concept}
\textbf{Multi-tape TM}: Multiple tapes, same power as single-tape TM

\textbf{Non-deterministic TM}: Multiple transitions possible, same power as deterministic TM

\textbf{Universal Turing Machine}: TM that can simulate any other TM

\textbf{All variants have same computational power}
\end{concept}

\section{Decidability}

\subsection{Decidable and Recognizable}
\begin{concept}
\textbf{Decidable Language} (Recursive):
\begin{itemize}
    \item TM halts and accepts strings in language
    \item TM halts and rejects strings not in language
    \item TM always halts
\end{itemize}

\textbf{Recognizable Language} (Recursively Enumerable):
\begin{itemize}
    \item TM accepts strings in language
    \item TM may reject or loop on strings not in language
\end{itemize}

\textbf{Relationships}:
\begin{itemize}
    \item Every decidable language is recognizable
    \item Not every recognizable language is decidable
\end{itemize}
\end{concept}

\subsection{Undecidable Problems}
\begin{concept}
\textbf{Halting Problem}:
\begin{itemize}
    \item Given TM $M$ and input $w$, does $M$ halt on $w$?
    \item Undecidable
    \item Proved using diagonalization
\end{itemize}

\textbf{Other Undecidable Problems}:
\begin{itemize}
    \item Does TM accept empty language?
    \item Are two CFGs equivalent?
    \item Is a CFG ambiguous?
    \item Post Correspondence Problem
\end{itemize}
\end{concept}

\subsection{Rice's Theorem}
\begin{theorem}
\textbf{Rice's Theorem}: Any non-trivial property of recognizable languages is undecidable.

\textbf{Non-trivial}: Some TMs have the property, some don't
\end{theorem}

\section{Complexity Theory}

\subsection{Time Complexity Classes}
\begin{concept}
\textbf{P (Polynomial Time)}:
\begin{itemize}
    \item Problems solvable in polynomial time by deterministic TM
    \item $P = \bigcup_{k} TIME(n^k)$
    \item "Efficiently solvable"
\end{itemize}

\textbf{NP (Non-deterministic Polynomial)}:
\begin{itemize}
    \item Problems verifiable in polynomial time
    \item Solvable in polynomial time by non-deterministic TM
    \item $P \subseteq NP$
\end{itemize}

\textbf{NP-Complete}:
\begin{itemize}
    \item In NP and NP-hard
    \item Hardest problems in NP
    \item If any NP-complete problem is in P, then $P = NP$
\end{itemize}

\textbf{NP-Hard}:
\begin{itemize}
    \item At least as hard as hardest NP problems
    \item Not necessarily in NP
\end{itemize}
\end{concept}

\subsection{P vs NP}
\begin{concept}
\textbf{P vs NP Question}: Does $P = NP$?
\begin{itemize}
    \item One of the most important open problems in CS
    \item Most believe $P \neq NP$
    \item Not yet proved
\end{itemize}
\end{concept}

\subsection{NP-Completeness}
\begin{concept}
\textbf{Showing NP-Completeness}:
\begin{enumerate}
    \item Show problem is in NP (polynomial verifier)
    \item Reduce a known NP-complete problem to it (polynomial reduction)
\end{enumerate}

\textbf{First NP-Complete Problem}: SAT (Boolean Satisfiability) - Cook-Levin Theorem

\textbf{Common NP-Complete Problems}:
\begin{itemize}
    \item SAT, 3-SAT
    \item Vertex Cover
    \item Hamiltonian Path/Cycle
    \item Traveling Salesman Problem (decision version)
    \item Subset Sum
    \item Graph Coloring
    \item Clique
\end{itemize}
\end{concept}

\chapter{Compiler Design}

\section{Phases of Compiler}

\subsection{Compiler Structure}
\begin{concept}
\textbf{Analysis Phase} (Front-end):
\begin{enumerate}
    \item \textbf{Lexical Analysis}: Source code $\to$ Tokens
    \item \textbf{Syntax Analysis}: Tokens $\to$ Parse Tree
    \item \textbf{Semantic Analysis}: Parse Tree $\to$ Annotated Tree
\end{enumerate}

\textbf{Synthesis Phase} (Back-end):
\begin{enumerate}
    \item \textbf{Intermediate Code Generation}
    \item \textbf{Code Optimization}
    \item \textbf{Code Generation}: Target code
\end{enumerate}

\textbf{Symbol Table}: Maintains information about identifiers

\textbf{Error Handler}: Reports errors at each phase
\end{concept}

\section{Lexical Analysis}

\subsection{Tokens and Lexemes}
\begin{concept}
\textbf{Lexical Analyzer (Scanner)}:
\begin{itemize}
    \item Reads source program character by character
    \item Groups into meaningful sequences (lexemes)
    \item Produces tokens
\end{itemize}

\textbf{Token}: <token-name, attribute-value>
\begin{itemize}
    \item Keywords: if, while, int
    \item Operators: +, -, *, /
    \item Identifiers: variable names
    \item Constants: numbers, strings
    \item Special symbols: ;, (, ), \{, \}
\end{itemize}

\textbf{Lexeme}: Sequence of characters matching token pattern
\end{concept}

\subsection{Regular Expressions for Tokens}
\begin{concept}
\textbf{Pattern Specification}:
\begin{itemize}
    \item Use regular expressions to specify token patterns
    \item Convert RE to NFA (Thompson's construction)
    \item Convert NFA to DFA (subset construction)
    \item Minimize DFA
    \item Implement as table-driven or direct-coded scanner
\end{itemize}
\end{concept}

\section{Syntax Analysis}

\subsection{Context-Free Grammars}
\begin{concept}
\textbf{Parser}: Checks if token stream conforms to grammar
\begin{itemize}
    \item Builds parse tree (or abstract syntax tree)
    \item Reports syntax errors
\end{itemize}

\textbf{Derivations}:
\begin{itemize}
    \item \textbf{Leftmost}: Expand leftmost non-terminal
    \item \textbf{Rightmost}: Expand rightmost non-terminal (canonical derivation)
\end{itemize}
\end{concept}

\subsection{Top-Down Parsing}
\begin{concept}
\textbf{Recursive Descent Parsing}:
\begin{itemize}
    \item Hand-coded parser
    \item One function per non-terminal
    \item May require backtracking
\end{itemize}

\textbf{LL(1) Parsing}:
\begin{itemize}
    \item Left-to-right scan
    \item Leftmost derivation
    \item 1 lookahead symbol
    \item Use predictive parsing table
    \item No backtracking
\end{itemize}

\textbf{FIRST and FOLLOW Sets}:
\begin{itemize}
    \item FIRST($\alpha$): Set of terminals beginning strings derived from $\alpha$
    \item FOLLOW(A): Set of terminals that can appear immediately after A
    \item Used to construct LL(1) parsing table
\end{itemize}
\end{concept}

\subsection{Bottom-Up Parsing}
\begin{concept}
\textbf{Shift-Reduce Parsing}:
\begin{itemize}
    \item Build parse tree from leaves to root
    \item \textbf{Shift}: Move input to stack
    \item \textbf{Reduce}: Replace stack top with non-terminal
    \item \textbf{Accept}: Input consumed, start symbol on stack
    \item \textbf{Error}: No action possible
\end{itemize}

\textbf{LR Parsing}:
\begin{itemize}
    \item Left-to-right scan
    \item Rightmost derivation (in reverse)
    \item More powerful than LL
\end{itemize}

\textbf{LR Variants}:
\begin{itemize}
    \item \textbf{SLR (Simple LR)}: Simple, less powerful
    \item \textbf{LALR (Look-Ahead LR)}: Used by yacc, good balance
    \item \textbf{CLR (Canonical LR)}: Most powerful, large tables
\end{itemize}
\end{concept}

\section{Semantic Analysis}

\subsection{Type Checking}
\begin{concept}
\textbf{Semantic Analyzer}:
\begin{itemize}
    \item Type checking
    \item Scope resolution
    \item Declaration checking
    \item Type conversions
\end{itemize}

\textbf{Attribute Grammars}:
\begin{itemize}
    \item \textbf{Synthesized Attributes}: Computed from children
    \item \textbf{Inherited Attributes}: Computed from parent/siblings
\end{itemize}
\end{concept}

\section{Intermediate Code Generation}

\subsection{Three-Address Code}
\begin{concept}
\textbf{Three-Address Code (TAC)}:
\begin{itemize}
    \item At most three operands per instruction
    \item Form: $x = y$ op $z$
    \item Easy to optimize
    \item Machine-independent
\end{itemize}

\textbf{Types of 3AC}:
\begin{itemize}
    \item Assignment: $x = y$ op $z$
    \item Unary: $x = $ op $y$
    \item Copy: $x = y$
    \item Unconditional jump: goto L
    \item Conditional jump: if $x$ relop $y$ goto L
    \item Procedure call: call p, n (n parameters)
\end{itemize}
\end{concept}

\section{Code Optimization}

\subsection{Optimization Techniques}
\begin{concept}
\textbf{Local Optimization} (within basic block):
\begin{itemize}
    \item Constant folding: Evaluate constants at compile time
    \item Constant propagation: Replace variables with constants
    \item Copy propagation: Replace copied variables
    \item Dead code elimination: Remove unreachable code
    \item Common subexpression elimination
\end{itemize}

\textbf{Global Optimization} (across basic blocks):
\begin{itemize}
    \item Loop optimization
    \item Code motion: Move loop-invariant code outside
    \item Induction variable elimination
    \item Strength reduction: Replace expensive operations
\end{itemize}
\end{concept}

\subsection{Data Flow Analysis}
\begin{concept}
\textbf{Reaching Definitions}:
\begin{itemize}
    \item Definition d reaches point p if path from d to p with no intervening definition
\end{itemize}

\textbf{Live Variable Analysis}:
\begin{itemize}
    \item Variable is live if its value may be used later
\end{itemize}

\textbf{Available Expressions}:
\begin{itemize}
    \item Expression available if computed and not modified
\end{itemize}
\end{concept}

\section{Code Generation}

\subsection{Target Code Generation}
\begin{concept}
\textbf{Code Generator}:
\begin{itemize}
    \item Translate intermediate code to target machine code
    \item Instruction selection
    \item Register allocation
    \item Instruction ordering
\end{itemize}

\textbf{Register Allocation}:
\begin{itemize}
    \item Graph coloring: Nodes = variables, Edges = interference
    \item k-colorable if k registers sufficient
\end{itemize}
\end{concept}

\chapter{Operating Systems}

\section{Process Management}

\subsection{Processes}
\begin{concept}
\textbf{Process}: Program in execution
\begin{itemize}
    \item \textbf{Process State}: New, Ready, Running, Waiting, Terminated
    \item \textbf{Process Control Block (PCB)}: Contains process state, PC, registers, memory limits, open files
\end{itemize}

\textbf{Process vs Program}:
\begin{itemize}
    \item Program: Passive entity (code on disk)
    \item Process: Active entity (program in execution)
\end{itemize}
\end{concept}

\subsection{Threads}
\begin{concept}
\textbf{Thread}: Lightweight process, basic unit of CPU utilization

\textbf{User Threads}:
\begin{itemize}
    \item Managed by user-level library
    \item Fast context switch
    \item OS unaware
\end{itemize}

\textbf{Kernel Threads}:
\begin{itemize}
    \item Managed by OS
    \item Slower context switch
    \item Better scheduling
\end{itemize}

\textbf{Benefits}:
\begin{itemize}
    \item Resource sharing
    \item Economy
    \item Scalability
    \item Responsiveness
\end{itemize}
\end{concept}

\subsection{CPU Scheduling}
\begin{concept}
\textbf{Scheduling Algorithms}:

\textbf{First-Come First-Served (FCFS)}:
\begin{itemize}
    \item Non-preemptive
    \item Simple but poor average waiting time
    \item Convoy effect
\end{itemize}

\textbf{Shortest Job First (SJF)}:
\begin{itemize}
    \item Optimal average waiting time
    \item Can be preemptive (SRTF) or non-preemptive
    \item Starvation possible
\end{itemize}

\textbf{Priority Scheduling}:
\begin{itemize}
    \item Each process has priority
    \item Can be preemptive or non-preemptive
    \item Starvation possible (solution: aging)
\end{itemize}

\textbf{Round Robin (RR)}:
\begin{itemize}
    \item Preemptive
    \item Time quantum
    \item Fair but context switching overhead
\end{itemize}

\textbf{Multilevel Queue}:
\begin{itemize}
    \item Separate queues for different priorities
    \item Each queue has own scheduling algorithm
\end{itemize}

\textbf{Multilevel Feedback Queue}:
\begin{itemize}
    \item Processes can move between queues
    \item Most general
\end{itemize}
\end{concept}

\begin{formula}
\textbf{Scheduling Metrics}:
\begin{itemize}
    \item \textbf{Turnaround Time}: Completion time - Arrival time
    \item \textbf{Waiting Time}: Turnaround time - Burst time
    \item \textbf{Response Time}: First response - Arrival time
    \item \textbf{CPU Utilization}: Percentage of time CPU is busy
    \item \textbf{Throughput}: Number of processes completed per unit time
\end{itemize}
\end{formula}

\section{Process Synchronization}

\subsection{Critical Section Problem}
\begin{concept}
\textbf{Critical Section}: Code accessing shared resources

\textbf{Requirements}:
\begin{enumerate}
    \item \textbf{Mutual Exclusion}: Only one process in critical section
    \item \textbf{Progress}: Selection of next process cannot be postponed indefinitely
    \item \textbf{Bounded Waiting}: Limit on waiting time
\end{enumerate}
\end{concept}

\subsection{Synchronization Mechanisms}
\begin{concept}
\textbf{Semaphore}:
\begin{itemize}
    \item Integer variable with atomic wait() and signal()
    \item \textbf{Binary Semaphore}: 0 or 1 (mutex)
    \item \textbf{Counting Semaphore}: Non-negative integer
\end{itemize}

\textbf{Mutex Lock}:
\begin{itemize}
    \item Binary lock (locked/unlocked)
    \item acquire() and release()
\end{itemize}

\textbf{Monitors}:
\begin{itemize}
    \item High-level synchronization construct
    \item Only one process active inside monitor
    \item Condition variables for waiting
\end{itemize}
\end{concept}

\subsection{Classic Synchronization Problems}
\begin{concept}
\textbf{Producer-Consumer Problem}:
\begin{itemize}
    \item Bounded buffer
    \item Semaphores: empty, full, mutex
\end{itemize}

\textbf{Readers-Writers Problem}:
\begin{itemize}
    \item Multiple readers or one writer
    \item Reader preference or Writer preference
\end{itemize}

\textbf{Dining Philosophers Problem}:
\begin{itemize}
    \item Five philosophers, five forks
    \item Deadlock possible
    \item Solutions: Asymmetric, resource hierarchy, waiter
\end{itemize}
\end{concept}

\subsection{Deadlocks}
\begin{concept}
\textbf{Deadlock}: Set of processes where each waits for resource held by another

\textbf{Necessary Conditions}:
\begin{enumerate}
    \item \textbf{Mutual Exclusion}: Non-sharable resources
    \item \textbf{Hold and Wait}: Process holds while waiting
    \item \textbf{No Preemption}: Resources not forcibly taken
    \item \textbf{Circular Wait}: Circular chain of waiting
\end{enumerate}

\textbf{Handling Deadlocks}:
\begin{itemize}
    \item \textbf{Prevention}: Negate one of four conditions
    \item \textbf{Avoidance}: Banker's algorithm (safe state)
    \item \textbf{Detection and Recovery}: Detect cycles, kill or preempt
    \item \textbf{Ignore}: Ostrich algorithm
\end{itemize}
\end{concept}

\section{Memory Management}

\subsection{Memory Allocation}
\begin{concept}
\textbf{Contiguous Allocation}:
\begin{itemize}
    \item \textbf{Fixed Partitioning}: Internal fragmentation
    \item \textbf{Variable Partitioning}: External fragmentation
\end{itemize}

\textbf{Allocation Strategies}:
\begin{itemize}
    \item \textbf{First Fit}: Allocate first hole large enough
    \item \textbf{Best Fit}: Allocate smallest hole large enough
    \item \textbf{Worst Fit}: Allocate largest hole
\end{itemize}

\textbf{Fragmentation}:
\begin{itemize}
    \item \textbf{Internal}: Allocated memory larger than needed
    \item \textbf{External}: Free memory scattered
    \item \textbf{Compaction}: Shuffle memory to collect free space
\end{itemize}
\end{concept}

\subsection{Paging}
\begin{concept}
\textbf{Paging}:
\begin{itemize}
    \item Divide physical memory into fixed-size frames
    \item Divide logical memory into same-size pages
    \item No external fragmentation
    \item Small internal fragmentation (last page)
\end{itemize}

\textbf{Address Translation}:
\begin{itemize}
    \item Logical address = (page number, offset)
    \item Page table maps page number to frame number
    \item Physical address = (frame number, offset)
\end{itemize}

\textbf{Page Table Structure}:
\begin{itemize}
    \item \textbf{Single-Level}: Simple but large
    \item \textbf{Hierarchical}: Multi-level page tables
    \item \textbf{Hashed}: Hash table
    \item \textbf{Inverted}: One entry per frame (not per page)
\end{itemize}
\end{concept}

\subsection{Segmentation}
\begin{concept}
\textbf{Segmentation}:
\begin{itemize}
    \item Divide memory into variable-size segments
    \item Logical view (code, data, stack)
    \item Segment table: base and limit for each segment
    \item External fragmentation possible
\end{itemize}

\textbf{Segmentation with Paging}:
\begin{itemize}
    \item Combine benefits of both
    \item Each segment is paged
\end{itemize}
\end{concept}

\subsection{Virtual Memory}
\begin{concept}
\textbf{Virtual Memory}:
\begin{itemize}
    \item Separation of logical and physical memory
    \item Only part of program needs to be in memory
    \item Allows very large programs
\end{itemize}

\textbf{Demand Paging}:
\begin{itemize}
    \item Load page only when needed
    \item \textbf{Page Fault}: Requested page not in memory
    \item Trap to OS, load page from disk
\end{itemize}

\textbf{Page Replacement}:
\begin{itemize}
    \item When memory full, replace a page
    \item \textbf{Reference String}: Sequence of page references
\end{itemize}
\end{concept}

\subsection{Page Replacement Algorithms}
\begin{concept}
\textbf{FIFO (First In First Out)}:
\begin{itemize}
    \item Replace oldest page
    \item Belady's anomaly possible
\end{itemize}

\textbf{Optimal}:
\begin{itemize}
    \item Replace page not used for longest future time
    \item Theoretical, used for comparison
\end{itemize}

\textbf{LRU (Least Recently Used)}:
\begin{itemize}
    \item Replace page not used for longest past time
    \item Good performance but expensive
    \item Implementation: Counter, Stack
\end{itemize}

\textbf{Second Chance (Clock)}:
\begin{itemize}
    \item FIFO with reference bit
    \item Give page second chance if referenced
\end{itemize}

\textbf{LFU (Least Frequently Used)}:
\begin{itemize}
    \item Replace page with lowest reference count
\end{itemize}
\end{concept}

\subsection{Thrashing}
\begin{concept}
\textbf{Thrashing}: Process spending more time paging than executing
\begin{itemize}
    \item Occurs when too many processes, insufficient frames
    \item CPU utilization drops
\end{itemize}

\textbf{Solutions}:
\begin{itemize}
    \item \textbf{Working Set Model}: Keep pages recently referenced
    \item \textbf{Page Fault Frequency}: Monitor page fault rate
    \item Reduce degree of multiprogramming
\end{itemize}
\end{concept}

\section{File Systems}

\subsection{File Concepts}
\begin{concept}
\textbf{File}: Named collection of related information
\begin{itemize}
    \item \textbf{Attributes}: Name, type, location, size, permissions, timestamps
    \item \textbf{Operations}: Create, open, read, write, seek, delete, truncate, close
\end{itemize}

\textbf{File Types}:
\begin{itemize}
    \item Regular files
    \item Directories
    \item Special files (devices)
\end{itemize}

\textbf{Access Methods}:
\begin{itemize}
    \item \textbf{Sequential}: Read/write in order
    \item \textbf{Direct/Random}: Access any location
    \item \textbf{Indexed}: Use index for access
\end{itemize}
\end{concept}

\subsection{Directory Structure}
\begin{concept}
\textbf{Directory Operations}:
\begin{itemize}
    \item Search, create, delete file
    \item List directory
    \item Rename, traverse file system
\end{itemize}

\textbf{Directory Structures}:
\begin{itemize}
    \item \textbf{Single-Level}: All files in one directory
    \item \textbf{Two-Level}: Separate directory per user
    \item \textbf{Tree-Structured}: Hierarchical
    \item \textbf{Acyclic Graph}: Shared subdirectories (hard/soft links)
    \item \textbf{General Graph}: Arbitrary structure (cycles possible)
\end{itemize}
\end{concept}

\subsection{File Allocation Methods}
\begin{concept}
\textbf{Contiguous Allocation}:
\begin{itemize}
    \item Files stored in contiguous blocks
    \item Fast sequential and direct access
    \item External fragmentation
\end{itemize}

\textbf{Linked Allocation}:
\begin{itemize}
    \item Each block points to next
    \item No external fragmentation
    \item Only sequential access
    \item Pointer overhead
\end{itemize}

\textbf{Indexed Allocation}:
\begin{itemize}
    \item Index block contains pointers to data blocks
    \item Direct access
    \item Index overhead
    \item Schemes: Linked, multilevel, combined (UNIX i-node)
\end{itemize}
\end{concept}

\subsection{Free Space Management}
\begin{concept}
\textbf{Free Space List Methods}:
\begin{itemize}
    \item \textbf{Bit Vector/Bitmap}: One bit per block
    \item \textbf{Linked List}: Free blocks linked
    \item \textbf{Grouping}: First free block contains addresses of n free blocks
    \item \textbf{Counting}: Store address and count of contiguous free blocks
\end{itemize}
\end{concept}

\section{I/O Systems}

\subsection{I/O Hardware}
\begin{concept}
\textbf{I/O Devices}:
\begin{itemize}
    \item \textbf{Block Devices}: Fixed-size blocks (disks)
    \item \textbf{Character Devices}: Stream of characters (keyboard, printer)
\end{itemize}

\textbf{Device Controllers}:
\begin{itemize}
    \item Electronics that operate device
    \item Has registers for communication with CPU
\end{itemize}
\end{concept}

\subsection{Disk Scheduling}
\begin{concept}
\textbf{Disk Access Time} = Seek time + Rotational latency + Transfer time

\textbf{Disk Scheduling Algorithms}:

\textbf{FCFS}: First Come First Served - No optimization

\textbf{SSTF (Shortest Seek Time First)}:
\begin{itemize}
    \item Select request with minimum seek time
    \item May cause starvation
\end{itemize}

\textbf{SCAN (Elevator)}:
\begin{itemize}
    \item Move in one direction, service requests
    \item Reverse at end
\end{itemize}

\textbf{C-SCAN (Circular SCAN)}:
\begin{itemize}
    \item Like SCAN but return to start without servicing
    \item More uniform wait time
\end{itemize}

\textbf{LOOK}:
\begin{itemize}
    \item Like SCAN but reverse at last request
\end{itemize}

\textbf{C-LOOK}:
\begin{itemize}
    \item Like C-SCAN but return to first request
\end{itemize}
\end{concept}

\chapter{Database Management Systems}

\section{Database Concepts}

\subsection{DBMS Basics}
\begin{concept}
\textbf{Database}: Organized collection of data

\textbf{DBMS Advantages}:
\begin{itemize}
    \item Data independence
    \item Efficient data access
    \item Data integrity and security
    \item Concurrent access
    \item Crash recovery
    \item Reduced redundancy
\end{itemize}

\textbf{Data Models}:
\begin{itemize}
    \item Relational
    \item Hierarchical
    \item Network
    \item Object-oriented
\end{itemize}
\end{concept}

\subsection{Three-Schema Architecture}
\begin{concept}
\begin{itemize}
    \item \textbf{External/View Level}: User views
    \item \textbf{Conceptual/Logical Level}: Logical structure
    \item \textbf{Internal/Physical Level}: Physical storage
\end{itemize}

\textbf{Data Independence}:
\begin{itemize}
    \item \textbf{Logical}: Change conceptual schema without changing views
    \item \textbf{Physical}: Change physical schema without changing conceptual
\end{itemize}
\end{concept}

\section{Relational Model}

\subsection{Relational Model Concepts}
\begin{concept}
\textbf{Relation}: Table with rows (tuples) and columns (attributes)
\begin{itemize}
    \item \textbf{Schema}: R(A1, A2, ..., An)
    \item \textbf{Domain}: Set of allowed values for attribute
    \item \textbf{Degree}: Number of attributes
    \item \textbf{Cardinality}: Number of tuples
\end{itemize}

\textbf{Keys}:
\begin{itemize}
    \item \textbf{Super Key}: Uniquely identifies tuple
    \item \textbf{Candidate Key}: Minimal super key
    \item \textbf{Primary Key}: Chosen candidate key
    \item \textbf{Foreign Key}: References primary key of another relation
\end{itemize}
\end{concept}

\subsection{Relational Algebra}
\begin{concept}
\textbf{Basic Operations}:
\begin{itemize}
    \item \textbf{Select} ($\sigma$): Select rows satisfying condition
    \item \textbf{Project} ($\pi$): Select columns
    \item \textbf{Union} ($\cup$): Tuples in either relation (compatible)
    \item \textbf{Set Difference} ($-$): Tuples in first but not second
    \item \textbf{Cartesian Product} ($\times$): All combinations
    \item \textbf{Rename} ($\rho$): Rename relation/attributes
\end{itemize}

\textbf{Derived Operations}:
\begin{itemize}
    \item \textbf{Intersection} ($\cap$): Tuples in both
    \item \textbf{Join} ($\bowtie$): Combine related tuples
    \item \textbf{Division} ($\div$): Tuples in R associated with all in S
\end{itemize}

\textbf{Join Types}:
\begin{itemize}
    \item \textbf{Natural Join}: Join on common attributes
    \item \textbf{Theta Join}: Join with condition
    \item \textbf{Equi Join}: Theta join with equality
    \item \textbf{Outer Join}: Include unmatched tuples (Left, Right, Full)
\end{itemize}
\end{concept}

\section{SQL}

\subsection{SQL Basics}
\begin{concept}
\textbf{Data Definition Language (DDL)}:
\begin{itemize}
    \item CREATE, ALTER, DROP, TRUNCATE
\end{itemize}

\textbf{Data Manipulation Language (DML)}:
\begin{itemize}
    \item SELECT, INSERT, UPDATE, DELETE
\end{itemize}

\textbf{Data Control Language (DCL)}:
\begin{itemize}
    \item GRANT, REVOKE
\end{itemize}

\textbf{Transaction Control}:
\begin{itemize}
    \item COMMIT, ROLLBACK, SAVEPOINT
\end{itemize}
\end{concept}

\subsection{Query Structure}
\begin{concept}
\textbf{Basic SELECT}:
\begin{verbatim}
SELECT [DISTINCT] attributes
FROM tables
WHERE condition
GROUP BY attributes
HAVING group_condition
ORDER BY attributes [ASC|DESC]
\end{verbatim}

\textbf{Aggregate Functions}:
\begin{itemize}
    \item COUNT, SUM, AVG, MAX, MIN
\end{itemize}

\textbf{Joins in SQL}:
\begin{itemize}
    \item INNER JOIN
    \item LEFT/RIGHT/FULL OUTER JOIN
    \item CROSS JOIN
    \item SELF JOIN
\end{itemize}

\textbf{Nested Queries}:
\begin{itemize}
    \item Subquery in WHERE clause
    \item IN, NOT IN, EXISTS, NOT EXISTS
    \item ANY, ALL operators
\end{itemize}
\end{concept}

\section{Normalization}

\subsection{Functional Dependencies}
\begin{concept}
\textbf{Functional Dependency}: $X \to Y$ means X determines Y
\begin{itemize}
    \item If two tuples agree on X, they must agree on Y
\end{itemize}

\textbf{Armstrong's Axioms}:
\begin{itemize}
    \item \textbf{Reflexivity}: If $Y \subseteq X$, then $X \to Y$
    \item \textbf{Augmentation}: If $X \to Y$, then $XZ \to YZ$
    \item \textbf{Transitivity}: If $X \to Y$ and $Y \to Z$, then $X \to Z$
\end{itemize}

\textbf{Derived Rules}:
\begin{itemize}
    \item Union, Decomposition, Composition, Pseudo-transitivity
\end{itemize}

\textbf{Closure}:
\begin{itemize}
    \item $X^+$: Set of all attributes functionally determined by X
\end{itemize}
\end{concept}

\subsection{Normal Forms}
\begin{concept}
\textbf{1NF (First Normal Form)}:
\begin{itemize}
    \item Atomic values (no multi-valued attributes)
    \item No repeating groups
\end{itemize}

\textbf{2NF}:
\begin{itemize}
    \item 1NF + No partial dependencies
    \item Non-prime attributes fully dependent on candidate key
\end{itemize}

\textbf{3NF}:
\begin{itemize}
    \item 2NF + No transitive dependencies
    \item For $X \to A$: Either X is super key OR A is prime attribute
\end{itemize}

\textbf{BCNF (Boyce-Codd NF)}:
\begin{itemize}
    \item For every $X \to A$: X must be super key
    \item Stricter than 3NF
\end{itemize}

\textbf{4NF}:
\begin{itemize}
    \item BCNF + No multi-valued dependencies
\end{itemize}
\end{concept}

\subsection{Decomposition}
\begin{concept}
\textbf{Lossless Join Decomposition}:
\begin{itemize}
    \item Can reconstruct original relation
    \item Natural join of decomposed relations = original
\end{itemize}

\textbf{Dependency Preservation}:
\begin{itemize}
    \item All FDs can be checked on individual relations
\end{itemize}

\textbf{Trade-off}:
\begin{itemize}
    \item BCNF guarantees lossless join
    \item 3NF guarantees both lossless join and dependency preservation
\end{itemize}
\end{concept}

\section{Transactions}

\subsection{Transaction Concepts}
\begin{concept}
\textbf{Transaction}: Unit of work (sequence of operations)

\textbf{ACID Properties}:
\begin{itemize}
    \item \textbf{Atomicity}: All or nothing
    \item \textbf{Consistency}: Database remains consistent
    \item \textbf{Isolation}: Concurrent transactions don't interfere
    \item \textbf{Durability}: Committed changes persist
\end{itemize}

\textbf{Transaction States}:
\begin{itemize}
    \item Active, Partially Committed, Committed, Failed, Aborted
\end{itemize}
\end{concept}

\subsection{Concurrency Control}
\begin{concept}
\textbf{Concurrency Problems}:
\begin{itemize}
    \item \textbf{Lost Update}: Two transactions update same data
    \item \textbf{Dirty Read}: Read uncommitted data
    \item \textbf{Non-repeatable Read}: Different values in same transaction
    \item \textbf{Phantom Read}: Different rows in same query
\end{itemize}

\textbf{Serializability}:
\begin{itemize}
    \item \textbf{Serial Schedule}: Transactions execute one after another
    \item \textbf{Serializable Schedule}: Equivalent to some serial schedule
    \item \textbf{Conflict Serializable}: Can transform to serial by swapping non-conflicting operations
    \item \textbf{View Serializable}: Same final state and reads as serial schedule
\end{itemize}
\end{concept}

\subsection{Locking}
\begin{concept}
\textbf{Lock Types}:
\begin{itemize}
    \item \textbf{Shared Lock (S)}: Read lock
    \item \textbf{Exclusive Lock (X)}: Write lock
\end{itemize}

\textbf{Two-Phase Locking (2PL)}:
\begin{itemize}
    \item \textbf{Growing Phase}: Acquire locks, no release
    \item \textbf{Shrinking Phase}: Release locks, no acquire
    \item Guarantees conflict serializability
\end{itemize}

\textbf{Variants}:
\begin{itemize}
    \item \textbf{Strict 2PL}: Hold exclusive locks until commit
    \item \textbf{Rigorous 2PL}: Hold all locks until commit
\end{itemize}

\textbf{Deadlock}:
\begin{itemize}
    \item Circular wait for locks
    \item Detection: Wait-for graph (cycle detection)
    \item Prevention: Timestamp ordering, timeout
\end{itemize}
\end{concept}

\subsection{Timestamp Ordering}
\begin{concept}
\textbf{Timestamp Protocol}:
\begin{itemize}
    \item Each transaction has timestamp
    \item Ensure serialization order matches timestamp order
    \item Read/Write timestamps for data items
\end{itemize}

\textbf{Thomas Write Rule}: Ignore obsolete writes

\textbf{Advantage}: No deadlocks

\textbf{Disadvantage}: Starvation possible, cascading rollbacks
\end{concept}

\section{Indexing}

\subsection{Index Structures}
\begin{concept}
\textbf{Ordered Indices}:
\begin{itemize}
    \item \textbf{Primary Index}: On primary key, sorted
    \item \textbf{Secondary Index}: On non-primary key
    \item \textbf{Clustering Index}: On non-key, sorted
\end{itemize}

\textbf{Dense vs Sparse}:
\begin{itemize}
    \item \textbf{Dense}: Index entry for every record
    \item \textbf{Sparse}: Index entry for some records
\end{itemize}

\textbf{Multi-level Index}: Index on index
\end{concept}

\subsection{B+ Trees}
\begin{concept}
\textbf{B+ Tree}: Most common index structure
\begin{itemize}
    \item Balanced tree
    \item All keys in leaves
    \item Leaves linked (sequential access)
    \item Fast search, insert, delete: $O(\log n)$
\end{itemize}

\textbf{Properties}:
\begin{itemize}
    \item All paths root to leaf have same length
    \item Internal nodes: at least $\lceil n/2 \rceil$ children (except root)
    \item Root: at least 2 children (unless leaf)
\end{itemize}
\end{concept}

\subsection{Hashing}
\begin{concept}
\textbf{Static Hashing}:
\begin{itemize}
    \item Fixed number of buckets
    \item Hash function maps key to bucket
    \item Fast for equality searches
\end{itemize}

\textbf{Dynamic Hashing}:
\begin{itemize}
    \item \textbf{Extendible Hashing}: Directory doubles when full
    \item \textbf{Linear Hashing}: Buckets added one at a time
\end{itemize}
\end{concept}

\chapter{Computer Networks}

\section{Network Fundamentals}

\subsection{OSI Model}
\begin{concept}
\textbf{7 Layers} (Bottom to Top):
\begin{enumerate}
    \item \textbf{Physical}: Bits on wire (cables, hubs)
    \item \textbf{Data Link}: Frame delivery (MAC, switches, Ethernet)
    \item \textbf{Network}: Packet routing (IP, routers)
    \item \textbf{Transport}: End-to-end communication (TCP, UDP)
    \item \textbf{Session}: Session management
    \item \textbf{Presentation}: Data format, encryption
    \item \textbf{Application}: User applications (HTTP, FTP, SMTP)
\end{enumerate}
\end{concept}

\subsection{TCP/IP Model}
\begin{concept}
\textbf{4 Layers}:
\begin{enumerate}
    \item \textbf{Network Access/Link}: Physical + Data Link
    \item \textbf{Internet}: IP layer
    \item \textbf{Transport}: TCP, UDP
    \item \textbf{Application}: HTTP, FTP, SMTP, DNS
\end{enumerate}
\end{concept}

\section{Data Link Layer}

\subsection{Framing and Error Detection}
\begin{concept}
\textbf{Framing Methods}:
\begin{itemize}
    \item Character/Byte stuffing
    \item Bit stuffing
\end{itemize}

\textbf{Error Detection}:
\begin{itemize}
    \item \textbf{Parity}: Single bit parity (even/odd)
    \item \textbf{Checksum}: Sum of data
    \item \textbf{CRC (Cyclic Redundancy Check)}: Polynomial division
        \begin{itemize}
            \item More reliable than checksum
            \item Can detect burst errors
        \end{itemize}
\end{itemize}

\textbf{Error Correction}:
\begin{itemize}
    \item \textbf{Hamming Code}: Can detect 2-bit errors, correct 1-bit error
\end{itemize}
\end{concept}

\subsection{Flow Control}
\begin{concept}
\textbf{Stop-and-Wait}:
\begin{itemize}
    \item Send one frame, wait for ACK
    \item Simple but inefficient
\end{itemize}

\textbf{Sliding Window}:
\begin{itemize}
    \item Send multiple frames before ACK
    \item Window size determines number
    \item \textbf{Go-Back-N}: Retransmit from error
    \item \textbf{Selective Repeat}: Retransmit only error frames
\end{itemize}
\end{concept}

\subsection{Multiple Access Protocols}
\begin{concept}
\textbf{ALOHA}:
\begin{itemize}
    \item Pure ALOHA: Transmit anytime, 18\% efficiency
    \item Slotted ALOHA: Time slots, 37\% efficiency
\end{itemize}

\textbf{CSMA (Carrier Sense Multiple Access)}:
\begin{itemize}
    \item Listen before transmit
    \item Persistent, Non-persistent, p-persistent
\end{itemize}

\textbf{CSMA/CD (Collision Detection)}:
\begin{itemize}
    \item Detect collision during transmission
    \item Stop, wait random time, retry
    \item Used in Ethernet
\end{itemize}

\textbf{CSMA/CA (Collision Avoidance)}:
\begin{itemize}
    \item Try to avoid collisions
    \item Used in WiFi (802.11)
\end{itemize}

\textbf{Token Ring/Bus}:
\begin{itemize}
    \item Token circulates
    \item Only token holder can transmit
    \item No collisions
\end{itemize}
\end{concept}

\subsection{Ethernet}
\begin{concept}
\textbf{Ethernet Frame}:
\begin{itemize}
    \item Preamble (7 bytes) + SFD (1 byte)
    \item Destination MAC (6) + Source MAC (6)
    \item Type/Length (2)
    \item Data (46-1500)
    \item CRC (4)
\end{itemize}

\textbf{MAC Address}: 48-bit physical address

\textbf{Switches}: Forward based on MAC address table
\end{concept}

\section{Network Layer}

\subsection{IP Addressing}
\begin{concept}
\textbf{IPv4}: 32-bit address (4 octets)
\begin{itemize}
    \item \textbf{Class A}: 0.0.0.0 - 127.255.255.255 (8-bit network)
    \item \textbf{Class B}: 128.0.0.0 - 191.255.255.255 (16-bit network)
    \item \textbf{Class C}: 192.0.0.0 - 223.255.255.255 (24-bit network)
    \item \textbf{Class D}: Multicast (224-239)
    \item \textbf{Class E}: Reserved (240-255)
\end{itemize}

\textbf{Private Addresses}:
\begin{itemize}
    \item 10.0.0.0/8
    \item 172.16.0.0/12
    \item 192.168.0.0/16
\end{itemize}

\textbf{Special Addresses}:
\begin{itemize}
    \item 127.0.0.1: Loopback
    \item 0.0.0.0: Default route
    \item 255.255.255.255: Broadcast
\end{itemize}
\end{concept}

\subsection{Subnetting}
\begin{concept}
\textbf{Subnet Mask}: Separates network and host portions

\textbf{CIDR (Classless Inter-Domain Routing)}:
\begin{itemize}
    \item Notation: IP/prefix-length (e.g., 192.168.1.0/24)
    \item Flexible network sizes
    \item Reduces routing table size
\end{itemize}

\textbf{Subnet Calculations}:
\begin{itemize}
    \item Network address: IP AND Subnet mask
    \item Broadcast address: Network OR (NOT Subnet mask)
    \item Number of hosts: $2^{(32-prefix)} - 2$
\end{itemize}
\end{concept}

\subsection{IPv6}
\begin{concept}
\textbf{IPv6}: 128-bit address (8 groups of 4 hex digits)
\begin{itemize}
    \item Larger address space
    \item Simplified header
    \item No fragmentation at routers
    \item Built-in security (IPSec)
    \item No broadcast (multicast instead)
\end{itemize}

\textbf{Address Types}:
\begin{itemize}
    \item Unicast
    \item Multicast
    \item Anycast
\end{itemize}
\end{concept}

\subsection{Routing Algorithms}
\begin{concept}
\textbf{Distance Vector}:
\begin{itemize}
    \item Bellman-Ford algorithm
    \item Each router knows distance to neighbors
    \item Share distance vectors
    \item Example: RIP (Routing Information Protocol)
    \item Count-to-infinity problem
\end{itemize}

\textbf{Link State}:
\begin{itemize}
    \item Dijkstra's algorithm
    \item Each router knows complete topology
    \item Flood link state information
    \item Example: OSPF (Open Shortest Path First)
    \item More memory and computation
\end{itemize}
\end{concept}

\subsection{Routing Protocols}
\begin{concept}
\textbf{Interior Gateway Protocols (IGP)}:
\begin{itemize}
    \item \textbf{RIP}: Distance vector, hop count metric, max 15 hops
    \item \textbf{OSPF}: Link state, cost metric, hierarchical
\end{itemize}

\textbf{Exterior Gateway Protocol (EGP)}:
\begin{itemize}
    \item \textbf{BGP (Border Gateway Protocol)}: Path vector, inter-AS routing
\end{itemize}
\end{concept}

\subsection{Network Address Translation (NAT)}
\begin{concept}
\textbf{NAT}:
\begin{itemize}
    \item Translates private to public IPs
    \item Conserves IPv4 addresses
    \item Provides basic security
    \item Breaks end-to-end principle
\end{itemize}

\textbf{Types}:
\begin{itemize}
    \item Static NAT: One-to-one
    \item Dynamic NAT: Pool of public IPs
    \item PAT (Port Address Translation): Many-to-one using ports
\end{itemize}
\end{concept}

\section{Transport Layer}

\subsection{UDP}
\begin{concept}
\textbf{UDP (User Datagram Protocol)}:
\begin{itemize}
    \item Connectionless
    \item Unreliable (no ACK, no retransmission)
    \item No flow control, no congestion control
    \item Low overhead (8-byte header)
    \item Fast
\end{itemize}

\textbf{UDP Header}:
\begin{itemize}
    \item Source Port (2 bytes)
    \item Destination Port (2 bytes)
    \item Length (2 bytes)
    \item Checksum (2 bytes)
\end{itemize}

\textbf{Use Cases}: DNS, DHCP, streaming, gaming
\end{concept}

\subsection{TCP}
\begin{concept}
\textbf{TCP (Transmission Control Protocol)}:
\begin{itemize}
    \item Connection-oriented
    \item Reliable (ACK, retransmission)
    \item Flow control (sliding window)
    \item Congestion control
    \item In-order delivery
    \item Full duplex
\end{itemize}

\textbf{TCP Header} (20-60 bytes):
\begin{itemize}
    \item Source/Destination Port (4 bytes)
    \item Sequence Number (4 bytes)
    \item Acknowledgment Number (4 bytes)
    \item Flags (SYN, ACK, FIN, RST, PSH, URG)
    \item Window Size (2 bytes)
    \item Checksum, Options
\end{itemize}
\end{concept}

\subsection{TCP Connection Management}
\begin{concept}
\textbf{Three-Way Handshake} (Connection Establishment):
\begin{enumerate}
    \item Client $\to$ Server: SYN
    \item Server $\to$ Client: SYN-ACK
    \item Client $\to$ Server: ACK
\end{enumerate}

\textbf{Connection Termination} (Four-Way Handshake):
\begin{enumerate}
    \item A $\to$ B: FIN
    \item B $\to$ A: ACK
    \item B $\to$ A: FIN
    \item A $\to$ B: ACK
\end{enumerate}
\end{concept}

\subsection{Flow and Congestion Control}
\begin{concept}
\textbf{Flow Control}:
\begin{itemize}
    \item Prevent sender from overwhelming receiver
    \item Receiver advertises window size
    \item Sliding window mechanism
\end{itemize}

\textbf{Congestion Control}:
\begin{itemize}
    \item Prevent network congestion
    \item \textbf{Slow Start}: Exponential increase
    \item \textbf{Congestion Avoidance}: Linear increase after threshold
    \item \textbf{Fast Retransmit}: Retransmit on 3 duplicate ACKs
    \item \textbf{Fast Recovery}: Halve window, continue
\end{itemize}
\end{concept}

\section{Application Layer}

\subsection{DNS}
\begin{concept}
\textbf{DNS (Domain Name System)}:
\begin{itemize}
    \item Hierarchical naming system
    \item Translates domain names to IP addresses
    \item Distributed database
\end{itemize}

\textbf{DNS Hierarchy}:
\begin{itemize}
    \item Root servers (.)
    \item TLD servers (.com, .org, .edu)
    \item Authoritative servers (example.com)
\end{itemize}

\textbf{Query Types}:
\begin{itemize}
    \item \textbf{Recursive}: DNS server resolves completely
    \item \textbf{Iterative}: DNS server returns referral
\end{itemize}

\textbf{Record Types}:
\begin{itemize}
    \item A: IPv4 address
    \item AAAA: IPv6 address
    \item MX: Mail server
    \item NS: Name server
    \item CNAME: Alias
\end{itemize}
\end{concept}

\subsection{HTTP}
\begin{concept}
\textbf{HTTP (HyperText Transfer Protocol)}:
\begin{itemize}
    \item Stateless protocol
    \item Request-response model
    \item Port 80 (HTTP), 443 (HTTPS)
\end{itemize}

\textbf{Request Methods}:
\begin{itemize}
    \item GET: Retrieve resource
    \item POST: Submit data
    \item PUT: Update resource
    \item DELETE: Delete resource
    \item HEAD: Get headers only
\end{itemize}

\textbf{Response Codes}:
\begin{itemize}
    \item 1xx: Informational
    \item 2xx: Success (200 OK)
    \item 3xx: Redirection (301 Moved, 304 Not Modified)
    \item 4xx: Client Error (404 Not Found, 403 Forbidden)
    \item 5xx: Server Error (500 Internal Error, 503 Unavailable)
\end{itemize}

\textbf{HTTP/2}:
\begin{itemize}
    \item Binary protocol
    \item Multiplexing
    \item Server push
    \item Header compression
\end{itemize}
\end{concept}

\subsection{Email Protocols}
\begin{concept}
\textbf{SMTP (Simple Mail Transfer Protocol)}:
\begin{itemize}
    \item Sending email
    \item Port 25, 587 (submission)
    \item Push protocol
\end{itemize}

\textbf{POP3 (Post Office Protocol)}:
\begin{itemize}
    \item Receiving email
    \item Port 110
    \item Download and delete from server
\end{itemize}

\textbf{IMAP (Internet Message Access Protocol)}:
\begin{itemize}
    \item Receiving email
    \item Port 143
    \item Keep messages on server
    \item Synchronization
\end{itemize}
\end{concept}

\subsection{DHCP}
\begin{concept}
\textbf{DHCP (Dynamic Host Configuration Protocol)}:
\begin{itemize}
    \item Automatic IP configuration
    \item Provides IP, subnet mask, gateway, DNS
\end{itemize}

\textbf{DORA Process}:
\begin{enumerate}
    \item \textbf{Discover}: Client broadcasts
    \item \textbf{Offer}: Server offers IP
    \item \textbf{Request}: Client requests IP
    \item \textbf{Acknowledge}: Server confirms
\end{enumerate}
\end{concept}

\section{Network Security}

\subsection{Cryptography Basics}
\begin{concept}
\textbf{Symmetric Key}:
\begin{itemize}
    \item Same key for encryption and decryption
    \item Fast
    \item Key distribution problem
    \item Examples: AES, DES, 3DES
\end{itemize}

\textbf{Asymmetric Key (Public Key)}:
\begin{itemize}
    \item Public key for encryption, private key for decryption
    \item Slower
    \item Solves key distribution
    \item Examples: RSA, ECC
\end{itemize}

\textbf{Hash Functions}:
\begin{itemize}
    \item One-way function
    \item Fixed output size
    \item Collision resistant
    \item Examples: SHA-256, MD5 (deprecated)
\end{itemize}
\end{concept}

\subsection{Security Protocols}
\begin{concept}
\textbf{SSL/TLS}:
\begin{itemize}
    \item Secure communication over network
    \item Uses certificates
    \item Handshake protocol
\end{itemize}

\textbf{IPSec}:
\begin{itemize}
    \item Security at network layer
    \item VPN implementation
    \item Authentication Header (AH)
    \item Encapsulating Security Payload (ESP)
\end{itemize}
\end{concept}

\end{document}
